% Theorem E: label-free local convergence under tangential excitation (PARADIGM.md, section 6).
\section{Theorem E: label-free local convergence under tangential excitation}\label{sec:thmE}

\paragraph{Assumptions.}
\begin{itemize}
  \item (S0) holds, in particular $k_e>0$.
  \item \textbf{Uniform tube.} There is $\rho>0$ such that $\varepsilon(x,t)=0$ whenever $\norm{x-x_m(t)}<\rho$, for
    all $t\ge0$.
  \item \textbf{Interior ideal.} $\mu^*$ is constant, with $\vartheta^*_j\in\operatorname{int}\ThKe$. Put
    $\delta:=\min_j\operatorname{dist}(\vartheta^*_j,\partial\ThKe)>0$.
  \item $x_m\in C^1$, and $x_m$ and $\dot x_m$ are bounded.
  \item Define
    $\Psi(t):=\bigl[\nabla_\vartheta\Phi(x_m(t);\vartheta^*_1),\dots,\nabla_\vartheta\Phi(x_m(t);\vartheta^*_m)\bigr]
    \in\R^{n\times mp}$.
  \item \textbf{(PE)} There are $T,\alpha>0$ with $\int_t^{t+T}\Psi^\T\Psi\dd r\succeq\alpha I_{mp}$ for all
    $t\ge0$.
\end{itemize}

(PE) forces $a^*_j\neq0$ and pairwise distinct pairs $(c^*_j,s^*_j)$.
\begin{itemize}
  \item If $a^*_j=0$, the $(c,s)$-columns of block $j$ vanish.
  \item If $(c^*_j,s^*_j)=(c^*_k,s^*_k)$, the kernel of $\int\Psi^\T\Psi\dd r$ contains the vector with $v$ in the
    $a$-slot of block $j$, $-v$ in the $a$-slot of block $k$, and zeros elsewhere.
\end{itemize}
In particular the $\vartheta^*_j$ are pairwise distinct, and
\begin{equation}\label{eq:thmE:gap}
  d^2:=\min_{\pi\neq\pi'}\frac1m\sum_i\norm{\vartheta^*_{\pi(i)}-\vartheta^*_{\pi'(i)}}^2_{G^{-1}}
  =\frac2m\cdot\min_{i\neq j}\norm{\vartheta^*_i-\vartheta^*_j}^2_{G^{-1}}>0,
\end{equation}
with $d^2:=+\infty$ for $m=1$. The closed form holds because a non-identity permutation moves at least two indices,
and the transposition of the $G^{-1}$-closest pair attains the bound.

\begin{thm}{Lemma}{E.0}[explicit strict Lyapunov function for the skew-symmetric mixed class]
Let $\Psi:[0,\infty)\to\R^{n\times q}$ be $C^1$ with $\norm{\Psi(t)}\le\Psi_M$ and $\norm{\dot\Psi(t)}\le\Psi_D$ for
all $t$, and persistently exciting: there are $T,\alpha>0$ with $\int_t^{t+T}\Psi(r)^\T\Psi(r)\dd r\succeq\alpha I_q$
for all $t\ge0$. Let $k>0$ and consider
\begin{equation}\label{eq:thmE:LTV}
  \dot e=-ke-\Psi(t)\zeta,\qquad \dot\zeta=\Psi(t)^\T e,\qquad w:=(e,\zeta)\in\R^{n+q}.
  \tag{LTV}
\end{equation}
Put $M:=\Psi^\T\Psi$ and
\begin{gather*}
  Q(t):=\int_t^{t+T}(t+T-r)M(r)\dd r,\qquad Q_M:=\frac{T^2\Psi_M^2}{2},\qquad
  c_1:=T^2\Psi_M^3+Tk\Psi_M+T\Psi_D,\\
  V(t,e,\zeta):=\tfrac12\norm{e}^2+\tfrac12\norm{\zeta}^2-\varepsilon\,\zeta^\T Q(t)\zeta+\varepsilon T\,e^\T\Psi(t)\zeta.
\end{gather*}
\begin{enumerate}[label=(\roman*)]
  \item \emph{Every admissible $\varepsilon$.} For every $\varepsilon>0$ with $\varepsilon(Q_M+T\Psi_M/2)<\frac12$,
    $V$ is $C^1$ and, along \textup{(LTV)},
    \begin{gather*}
      p_{\mathrm{lo}}\norm{w}^2\le V\le p_{\mathrm{hi}}\norm{w}^2,\qquad
      p_{\mathrm{lo}/\mathrm{hi}}:=\tfrac12\mp\varepsilon(Q_M+T\Psi_M/2),\\
      \dot V\le-(k-\varepsilon T\Psi_M^2)\norm{e}^2-\varepsilon\alpha\norm{\zeta}^2+\varepsilon c_1\norm{e}\norm{\zeta}.
    \end{gather*}
    Hence $\dot V\le-\gamma(\varepsilon)\norm{w}^2$ with
    $\gamma(\varepsilon):=\lambda_{\min}\bigl[\begin{smallmatrix}a&-c/2\\-c/2&b\end{smallmatrix}\bigr]$,
    $a:=k-\varepsilon T\Psi_M^2$, $b:=\varepsilon\alpha$, $c:=\varepsilon c_1$, whenever this is positive, and then
    $\norm{w(t)}\le(p_{\mathrm{hi}}/p_{\mathrm{lo}})^{1/2}e^{-(\gamma/(2p_{\mathrm{hi}}))(t-t_0)}\norm{w(t_0)}$ for
    all $t\ge t_0\ge0$.
  \item \emph{Closed form.} With
    $\varepsilon:=\min\bigl\{1/(4(Q_M+T\Psi_M/2)),\,k/(4T\Psi_M^2),\,k\alpha/c_1^2\bigr\}$, one has
    $p_{\mathrm{lo}}\ge\frac14$, $p_{\mathrm{hi}}\le\frac34$ and $\dot V\le-\gamma\norm{w}^2$ with
    $\gamma:=\min\{k/4,\,\varepsilon\alpha/2\}$. Hence
    \[
      \norm{w(t)}\le\sqrt3\cdot e^{-(2\gamma/3)(t-t_0)}\cdot\norm{w(t_0)}\qquad\text{for all }t\ge t_0\ge0.
    \]
  \item \emph{Dependence on $\Psi_D$ and on time rescaling.} The rate tends to $0$ as $\Psi_D\to\infty$, and it must:
    $\Psi=\sin t^2$ is bounded and PE with $\Psi_D=\infty$, and \textup{(LTV)} is then not uniformly exponentially
    stable (part~(f) of the proof). Under the time rescaling $\Psi_\omega(t):=\Psi(\omega t)$, with $\Psi_D>0$, the
    data transform as $(\alpha,T,\Psi_M,\Psi_D)\mapsto(\alpha/\omega,T/\omega,\Psi_M,\omega\Psi_D)$, so
    $c_1=T\Psi_D\cdot(1+O(1/\omega))$ and the closed-form rate is
    $k\alpha^2/(3\omega^2T^2\Psi_D^2)\cdot(1+O(1/\omega))$ as $\omega\to\infty$. This is the $\omega^{-2}$ law of
    Remark E.1.
\end{enumerate}
\end{thm}

\begin{proof}
\emph{(a) The Gramian weight.} $0\preceq Q(t)\preceq\int_t^{t+T}(t+T-r)\Psi_M^2\dd r\cdot I=Q_MI$. Since $M$ is
continuous, $Q$ is $C^1$ with $\dot Q(t)=\int_t^{t+T}M(r)\dd r-TM(t)$: the upper limit contributes
$(t+T-(t+T))M(t+T)=0$, the lower limit contributes $-TM(t)$, and differentiating the weight under the integral gives
$\int_t^{t+T}M$.

\noindent\emph{(b) Bounds on $V$.} $\abs{\zeta^\T Q\zeta}\le Q_M\norm{\zeta}^2$ and
$\abs{e^\T\Psi\zeta}\le\Psi_M\norm{e}\norm{\zeta}\le(\Psi_M/2)(\norm{e}^2+\norm{\zeta}^2)$.

\noindent\emph{(c) Three derivatives along \textup{(LTV)}.}
\begin{align*}
  \frac{\mathrm{d}}{\mathrm{d}t}\tfrac12(\norm{e}^2+\norm{\zeta}^2)
    &=-k\norm{e}^2-e^\T\Psi\zeta+\zeta^\T\Psi^\T e=-k\norm{e}^2,\\
  \frac{\mathrm{d}}{\mathrm{d}t}(-\zeta^\T Q\zeta)&=-\zeta^\T\dot Q\zeta-2\zeta^\T Q\dot\zeta
    =-\zeta^\T\Bigl(\int_t^{t+T}M\Bigr)\zeta+T\norm{\Psi\zeta}^2-2\zeta^\T Q\Psi^\T e\\
    &\le-\alpha\norm{\zeta}^2+T\norm{\Psi\zeta}^2-2\zeta^\T Q\Psi^\T e,\\
  \frac{\mathrm{d}}{\mathrm{d}t}(e^\T\Psi\zeta)&=\dot e^\T\Psi\zeta+e^\T\dot\Psi\zeta+e^\T\Psi\dot\zeta
    =-k\,e^\T\Psi\zeta-\norm{\Psi\zeta}^2+e^\T\dot\Psi\zeta+\norm{\Psi^\T e}^2.
\end{align*}
Adding them with weights $1$, $\varepsilon$ and $\varepsilon T$, the $\norm{\Psi\zeta}^2$ terms cancel:
\[
  \dot V\le-k\norm{e}^2-\varepsilon\alpha\norm{\zeta}^2-2\varepsilon\zeta^\T Q\Psi^\T e
  -\varepsilon Tk\,e^\T\Psi\zeta+\varepsilon T\,e^\T\dot\Psi\zeta+\varepsilon T\norm{\Psi^\T e}^2.
\]
The last four terms are at most $2\varepsilon Q_M\Psi_M\norm{e}\norm{\zeta}$,
$\varepsilon Tk\Psi_M\norm{e}\norm{\zeta}$, $\varepsilon T\Psi_D\norm{e}\norm{\zeta}$ and
$\varepsilon T\Psi_M^2\norm{e}^2$. This is (i); the quadratic form
$-a\norm{e}^2-b\norm{\zeta}^2+c\norm{e}\norm{\zeta}$ is at most
\[
  -\lambda_{\min}\begin{bmatrix}a&-c/2\\-c/2&b\end{bmatrix}\cdot\norm{w}^2,
\]
and $V\le p_{\mathrm{hi}}\norm{w}^2$, $\norm{w}^2\le V/p_{\mathrm{lo}}$ give the estimate.

\noindent\emph{(d) The closed form.} Young's inequality gives
$\varepsilon c_1\norm{e}\norm{\zeta}\le(\varepsilon\alpha/2)\norm{\zeta}^2+(\varepsilon c_1^2/(2\alpha))\norm{e}^2$.
With $\varepsilon\le k\alpha/c_1^2$ the second coefficient is at most $k/2$, and with
$\varepsilon\le k/(4T\Psi_M^2)$ the coefficient of $\norm{e}^2$ in (i) is at least $k-k/4$. So
$\dot V\le-(k/4)\norm{e}^2-(\varepsilon\alpha/2)\norm{\zeta}^2\le-\gamma\norm{w}^2$. Since
$V\le(3/4)\norm{w}^2$, $\dot V\le-(4\gamma/3)V$, and $\norm{w}^2\le4V$ gives
\[
  \norm{w(t)}^2\le4V(t)\le4e^{-(4\gamma/3)(t-t_0)}V(t_0)\le3e^{-(4\gamma/3)(t-t_0)}\norm{w(t_0)}^2.
\]

\noindent\emph{(e) Rescaling.} $\Psi_\omega$ has sup norm $\Psi_M$, derivative bound $\omega\Psi_D$, and
\[
  \int_t^{t+T/\omega}\Psi_\omega^\T\Psi_\omega\dd r=\frac1\omega\int_{\omega t}^{\omega t+T}\Psi^\T\Psi\dd r
  \succeq\frac\alpha\omega I.
\]
Substitute into $c_1$ and the closed form.

\noindent\emph{(f) $\Psi=\sin t^2$ is not uniformly exponentially stabilising} ($n=q=1$, $k>0$). Let
$F(s):=\int_s^\infty\sin r^2\dd r$, so that $\sin s^2=-F'(s)$ and $\abs{F(s)}\le C/s$ (Fresnel). From
$w(t_0)=(0,\zeta_0)$, $e(t)=-\int_{t_0}^te^{-k(t-s)}\sin(s^2)\zeta(s)\dd s$; integrating by parts and using
$\abs{\zeta}\le\norm{w(t_0)}$ (the energy $\frac12\norm{w}^2$ is non-increasing) and
$\abs{\dot\zeta}=\abs{\Psi e}\le\abs{e}$,
\[
  \abs{e(t)}\le\frac{C}{t_0}(2+1)\norm{w(t_0)}+\frac{C}{t_0}\int_{t_0}^te^{-k(t-s)}\abs{e(s)}\dd s,
\]
so $\sup_{t\ge t_0}\abs{e}\le3C\norm{w(t_0)}/\bigl(t_0(1-C/(kt_0))\bigr)$ for $t_0>C/k$. Hence
\[
  \tfrac12\norm{w(t_0+\Delta)}^2\ge\tfrac12\norm{w(t_0)}^2-k\Delta\Bigl(\frac{3C}{t_0}\Bigr)^2\norm{w(t_0)}^2(1+o(1)):
\]
the decay over any fixed $\Delta$ tends to $0$ as $t_0\to\infty$, which excludes a uniform exponential rate. The
three ratios quoted in step~3 of the proof of Theorem E ($0.9715$, $0.9942$, $0.9980$ for $t_0=5$, $20$, $50$ and
$\Delta=200$) match $1-\frac18\bigl(1/t_0-1/(t_0+200)\bigr)$ to three digits.
\end{proof}
\checkedby{\code{tests/test_excitation.py}: the two bounds on $V$ and $\dot V\le-\gamma\norm{w}^2$ at random $(t,w)$
for four regressors of different shapes and gains, with the closed-form $\varepsilon$ and with the $\varepsilon$
maximising the rate (\code{excitation.optimise_eps}; every admissible $\varepsilon$ gives a valid Lyapunov function,
so the maximum over the grid is a valid bound for the given PE data; the data $\Psi_M$, $\Psi_D$, $\alpha$ are
themselves grid maxima and a refinement-checked quadrature---the PE level over the $20\pi$ window is computed with
$800$ Gauss--Legendre nodes, and \code{excitation.pe_level} raises if a $2\times$ refinement disagrees---so the
ordering against the exact rate is numerical); the envelope along trajectories; the ordering explicit rate $\le$
exact Floquet rate; the $\omega^{-2}$ law; the slow-adaptation limit (a no-growth check of the ratio, not a
saturation check); and the explicit nonlinear envelope of Theorem E on the implemented closed loop. The round-5
referee verified the pointwise form $\lambda_{\max}(N(t))\le-\gamma$ of $\dot V\le-\gamma\norm{w}^2$ over all $w$,
with an independent quadrature, for ten regressor families ($q>n$, $q<n$, $k$ from $10^{-3}$ to $10^4$,
$\Psi_D/\Psi_M$ up to $100$, non-periodic PE, near-degenerate PE) and eight choices of $\varepsilon$, without finding
a violation.}

\begin{rmk}{Remark}{E.0a}[how conservative the explicit rate is]
The construction is of the Mazenc--Malisoff type, a backward--forward integral of the PE Gramian as a cross term; its
content here is the explicit, self-contained statement for the mixed class (LTV), which removes the dependence of
Theorem E on~\cite{S8}. The price is conservativeness, which the table of \code{exp_explicit_rate.py} $\to$
\code{results/explicit_rate/table.md} quantifies against the exact Floquet rate of the linearised loop.
\begin{itemize}
  \item \emph{Slow adaptation} ($T\Psi_M^2\ll k$). To leading order in the gains, $c_1\approx T(k\Psi_M+\Psi_D)$ and
    $\alpha=T\lambda_{\min}\langle\Psi^\T\Psi\rangle$ for a window equal to the period, so the closed-form bound
    $\frac23\gamma_{\mathrm{cf}}$ of (ii) is $k\alpha^2/(3c_1^2)$, and the rate with the optimal $\varepsilon$ of the
    same quadratic form is $k\alpha^2/c_1^2$ (maximise $\varepsilon\alpha-\varepsilon^2c_1^2/(4k)$ over
    $\varepsilon$). The exact rate is $\lambda_{\min}\langle\Psi^\T\Psi\rangle/k$ when $k$ is also large compared with
    the bandwidth of the reference (Remark E.1). Hence, in that doubly asymptotic regime,
    \[
      \frac{\text{exact}}{\text{explicit}}\approx3\kappa_\Psi\Bigl(1+\frac{\Psi_D}{k\Psi_M}\Bigr)^2\ \text{for (ii)},
      \qquad
      \approx\kappa_\Psi\Bigl(1+\frac{\Psi_D}{k\Psi_M}\Bigr)^2\ \text{for the optimised }\varepsilon,
    \]
    with $\kappa_\Psi:=\Psi_M^2/\lambda_{\min}\langle\Psi^\T\Psi\rangle$ the condition number of the averaged
    Gramian. The exact rate pays for ill-conditioning once, through $\lambda_{\min}$; the bound pays twice, because
    the Young step that absorbs the cross term $\varepsilon c_1\norm{e}\norm{\zeta}$ compares $c_1\propto\Psi_M$ with
    $\alpha\propto\lambda_{\min}$ (Lemma E.0$'$ below removes this by slaving the error to the slow manifold before
    the Young step). On the ideal of the Theorem E tests ($m=3$, $n=2$, $k_e=6$), driven by a two-tone
    reference of period $20\pi$, with $G=g\cdot\operatorname{diag}(2,2,1)$: $\kappa_\Psi=1.27\cdot10^5$,
    $(1+\Psi_D/(k_e\Psi_M))^2=3.84$, so the two formula values are $1.47\cdot10^6$ and $4.9\cdot10^5$. The measured
    ratio with the optimised $\varepsilon$ falls from $6.56\cdot10^{10}$ at $g=30$ to $4.22\cdot10^5$ at $g=0.01$
    ($5.63\cdot10^5$ with the closed-form $\varepsilon$, which the table evaluates with the exact $2\times2$ form and
    the exact $p_{\mathrm{hi}}$ rather than the bound (ii)). Two corrections account for the difference from
    $4.9\cdot10^5$: on this instance the exact rate is $0.66$ of $\lambda_{\min}\langle\Psi^\T\Psi\rangle/k$ at every
    $g$ ($k_e=6$ is not large against the reference's bandwidth; the two last columns of the table), and the
    $T^2\Psi_M^3$ term is still $13\%$ of $c_1$ at $g=0.01$ ($(c_1/c_1^\infty)^2=1.31$); their product,
    $0.66\times1.31=0.86$, brings $4.9\cdot10^5$ to $4.2\cdot10^5$. The ratio is still decreasing at $g=0.01$ and
    tends to about $0.66\cdot\kappa_\Psi(1+\Psi_D/(k_e\Psi_M))^2\approx3.2\cdot10^5$ as $g\to0$. In the mean-field
    normalisation $\Psi_G\propto m^{-1/2}$, so large networks sit in this regime at fixed geometry.
  \item \emph{Fast adaptation} ($T\Psi_M^2\gtrsim k$). The ratio grows further with the excitation strength, through
    the $T^2\Psi_M^3$ term of $c_1$ and the $Q_M$ term of $\varepsilon$.
  \item \emph{Remark E.1's instance} ($m=n=1$, $G=I$, $k_e=3$, $x_m=1.5\sin\omega t$). The ratio exact/explicit
    decreases from $954$ at $\omega=1$ to $17.6$ at $\omega=64$. Over $\omega\ge4$ the log-log slopes are $-1.97$
    (exact) and $-1.56$ (explicit); over the last octave, $\omega=32\to64$, the explicit slope is $-1.92$, approaching
    the asymptotic $\omega^{-2}$ of (iii).
  \item \emph{The window.} Any window length $T$ is admissible with its own $\alpha(T)$. On the test ideal the best
    window is half a period.
  \item \emph{Answered below.} Lemma E.0$'$ is a strict Lyapunov function whose rate is
    $\lambda_{\min}\langle\Psi^\T\Psi\rangle/k$ up to a factor $1-\delta$, above a threshold in $k$ set by the peak
    factors of the regressor: in the slow regime it pays for the ill-conditioning once.
\end{itemize}
\end{rmk}

\begin{thm}{Lemma}{E.0$'$}[the slow-adaptation regime: a strict Lyapunov function whose rate is the averaged rate]
Let $\Psi$ be $T$-periodic and $C^2$, $M:=\Psi^\T\Psi$, $\bar M:=\frac1T\int_0^TM\dd t\succ0$,
$\norm{v}_{\bar M}:=(v^\T\bar Mv)^{1/2}$, $\lambda_{\min}:=\lambda_{\min}(\bar M)$, and
\begin{gather*}
  \Psi_M:=\sup\norm{\Psi},\qquad \Psi_D:=\sup\norm{\dot\Psi},\qquad \Psi_{DD}:=\sup\norm{\ddot\Psi},\\
  \rho:=\sup_{t,\,v\neq0}\frac{\norm{\Psi(t)v}^2}{\norm{v}^2_{\bar M}},\qquad
  \rho_D:=\sup\frac{\norm{\dot\Psi(t)v}^2}{\norm{v}^2_{\bar M}},\qquad
  \rho_{DD}:=\sup\frac{\norm{\ddot\Psi(t)v}^2}{\norm{v}^2_{\bar M}},
\end{gather*}
the \emph{peak factors}: the largest ratio, over directions, between the instantaneous excitation of a direction (by
$\Psi$, $\dot\Psi$, $\ddot\Psi$) and its average excitation; $\rho\ge1$, and all three are finite because
$\bar M\succ0$. Write $\varepsilon:=1/k$ and let
\begin{align*}
  P_1(t)&:=-\frac2T\int_t^{t+T}(t+T-r)\bigl(M(r)-\bar M\bigr)\dd r,\qquad F:=P_1M+MP_1,\\
  S&:=\text{the zero-mean antiderivative of }F,\qquad P_2:=-(M-\bar M)+S,\\
  e_2(\zeta,t)&:=-\varepsilon\Psi\zeta+\varepsilon^2\dot\Psi\zeta-\varepsilon^3\Psi M\zeta,\qquad w:=e-e_2(\zeta,t),\\
  V(e,\zeta,t)&:=\tfrac12\zeta^\T\bigl(I+\varepsilon P_1(t)+\varepsilon^2P_2(t)\bigr)\zeta+\frac\beta2\norm{w}^2,
\end{align*}
where $P_1$ is the zero-mean antiderivative of $2(M-\bar M)$ and $e_2$ is the slow manifold of the error, expanded
through order $\varepsilon^3$. $F$ has zero mean: $\langle P_1M+MP_1\rangle=\langle P_1\rangle\bar M+\bar M\langle P_1\rangle
+\frac12\bigl\langle\frac{\mathrm{d}}{\mathrm{d}t}(P_1^2)\bigr\rangle=0$, so $S$ exists and the second-order averaged
correction $N_2:=\frac12\langle F\rangle$ vanishes identically. Put
\begin{align*}
  \pi_1&:=\tfrac T2\Psi_M(\sqrt\rho+1),\qquad
  \pi_2:=\Psi_M(\sqrt\rho+1)+\tfrac{T^2}8\Psi_M^3(2\sqrt\rho+1),\\
  p_1&:=\tfrac T2\Psi_M^2,\qquad p_2:=\Psi_M^2+\tfrac{T^2}4\Psi_M^4,\\
  c_1&:=\sqrt\rho+\varepsilon\Psi_M\pi_1+\varepsilon^2\Psi_M\pi_2,\\
  c_2&:=\varepsilon^2\sqrt{\rho_{DD}}+\varepsilon^3\bigl(3\Psi_D\Psi_M\sqrt\rho+2\Psi_M^2\sqrt{\rho_D}\bigr)
    +\varepsilon^4\bigl(2\Psi_M^4\sqrt\rho+\Psi_D\Psi_M\sqrt{\rho_D}\bigr)\\
  &\qquad+\varepsilon^5\bigl(\Psi_M^4\sqrt{\rho_D}+\Psi_D\Psi_M^3\sqrt\rho\bigr)+\varepsilon^6\Psi_M^6\sqrt\rho,\\
  d_w&:=\varepsilon\Psi_M^2+\varepsilon^2\Psi_D\Psi_M+\varepsilon^3\Psi_M^4,\\
  r_\zeta&:=\varepsilon^3\bigl(\pi_2\Psi_M\sqrt\rho+\pi_1\Psi_M\sqrt{\rho_D}+\Psi_M^2\rho\bigr)
    +\varepsilon^4\bigl(\pi_2\Psi_M\sqrt{\rho_D}+\pi_1\Psi_M^3\sqrt\rho\bigr)+\varepsilon^5\pi_2\Psi_M^3\sqrt\rho,\\
  \beta&:=2c_1^2/\delta,\qquad c_\zeta:=1-\delta-r_\zeta/\varepsilon.
\end{align*}
Fix $\delta\in(0,1)$. If
\[
  \text{(A)}\ c_1c_2\le\delta/2,\qquad \text{(B)}\ d_w\le k/4,\qquad \text{(C)}\ c_\zeta>0,\qquad
  \text{(D)}\ \varepsilon p_1+\varepsilon^2p_2<1,
\]
then along \textup{(LTV)}, for all $(e,\zeta,t)$,
\[
  \dot V\le-\varepsilon c_\zeta\,\zeta^\T\bar M\zeta-\frac{\beta k}4\norm{w}^2\le-\gamma V,\qquad
  \gamma:=\min\Bigl\{\frac{2c_\zeta\lambda_{\min}}{k(1+\varepsilon p_1+\varepsilon^2p_2)},\ \frac k2\Bigr\},
\]
and $\norm{(e,\zeta)(t)}\le Ce^{-(\gamma/2)(t-t_0)}\norm{(e,\zeta)(t_0)}$ with
\[
  C:=\bigl(1+\varepsilon\Psi_M+\varepsilon^2\Psi_D+\varepsilon^3\Psi_M^3\bigr)^2\cdot
  \Bigl(\frac{\max\{1+\varepsilon p_1+\varepsilon^2p_2,\ \beta\}}{\min\{1-\varepsilon p_1-\varepsilon^2p_2,\ \beta\}}\Bigr)^{1/2}.
\]

In the slow-adaptation limit ($k\to\infty$ at fixed $\Psi$) the conditions hold for every $\delta\in(0,1)$ once
$k\ge k_0(\delta)$, with $k_0\approx(4\rho\rho_{DD}/\delta^2)^{1/4}$ when (A) binds (small $\delta$, or
$T\Psi_M^2\lesssim1$; otherwise (C) sets $k_0$ through $r_\zeta$, as on the sinusoidal regressor below), and the rate
is $\gamma/2=(1-\delta)(1-\varepsilon p_1+O(k^{-2}))\lambda_{\min}/k$: the averaged rate
$\lambda_{\min}\langle\Psi^\T\Psi\rangle/k$ of Remark E.1 up to the factor $1-\delta$ and the first-order deficit
$\varepsilon p_1=T\Psi_M^2/(2k)$ of the quadratic form. The bound pays for the ill-conditioning of $\bar M$ once,
through $\lambda_{\min}$, as the exact rate does; the peak factors enter only the threshold $k_0$ and the overshoot
$C=O((\rho/\delta)^{1/2})$, not the rate.
\end{thm}
\status{Proved (below); checked pointwise by a certificate---the supremum of $\dot V/V$ over all states and over a
time grid is a generalised eigenvalue problem---on the sinusoidal regressor and on the slow-adaptation ideal of the
Theorem E tests, where the guaranteed rate is within 12\% ($\delta=0.1$) and 26\% ($\delta=0.2$) of the exact rate at
$k=400$ while Lemma E.0's is five orders of magnitude below it (\code{tests/test_excitation.py},
\code{exp_explicit_rate.py}, Table~\ref{tab:thmE:slowregime}). The round-6 referee re-derived every step and
certified the inequality independently on five regressor families without a failure
(\code{review/REVIEW_round6_E0prime.md}).}

\begin{proof}
Write $M$, $\Psi$, \dots\ for their values at $t$. Two facts about zero-mean antiderivatives are used throughout: if
$f$ is $T$-periodic with zero mean and $g$ is its zero-mean antiderivative, then
$\norm{g}_\infty\le(T/4)\norm{f}_\infty$---for the maximiser $t^*$ of $g$,
$g(t^*)=\frac1T\int_0^T(g(t^*)-g(s))\dd s\le\frac1T\int_0^T\operatorname{dist}(t^*,s)\norm{f}_\infty\dd s
=(T/4)\norm{f}_\infty$ with $\operatorname{dist}$ the distance on the circle, and the same for $-g$---and the same
bound holds for every bilinear coordinate $u^\T gv$, hence $\norm{g(t)v}\le(T/4)\sup_r\norm{f(r)v}$. As
$\norm{M-\bar M}\le\Psi_M^2$ (two positive semidefinite matrices of norm $\le\Psi_M^2$), this gives
$\norm{P_1}\le T\Psi_M^2/2=p_1$ and $\norm{P_1v}\le(T/2)\sup\norm{(M-\bar M)v}$, and $\norm{S}\le(T/4)\norm{F}_\infty$,
$\norm{Sv}\le(T/4)\sup\norm{Fv}$.

\noindent\emph{(a) The fast variable.} From $\dot e=-ke-\Psi\zeta$ and
$ke_2=-\Psi\zeta+\varepsilon\dot\Psi\zeta-\varepsilon^2\Psi M\zeta$,
\[
  \dot e=-kw-\varepsilon\dot\Psi\zeta+\varepsilon^2\Psi M\zeta.
\]
Differentiating $e_2$ along the flow, with
$\dot\zeta=\Psi^\T e=\Psi^\T w-\varepsilon M\zeta+\varepsilon^2\Psi^\T\dot\Psi\zeta-\varepsilon^3M^2\zeta$ and
$\dot M=\dot\Psi^\T\Psi+\Psi^\T\dot\Psi$, the terms of order $\varepsilon$ and $\varepsilon^2$ in $\zeta$ cancel and
\begin{align*}
  \dot w&=-kw+\varepsilon\Psi\Psi^\T w-\varepsilon^2\dot\Psi\Psi^\T w+\varepsilon^3\Psi M\Psi^\T w
    -\varepsilon^2\ddot\Psi\zeta+\varepsilon^3\bigl(2\dot\Psi M+\Psi\dot\Psi^\T\Psi+2\Psi\Psi^\T\dot\Psi\bigr)\zeta\\
  &\qquad-\varepsilon^4\bigl(2\Psi M^2+\dot\Psi\Psi^\T\dot\Psi\bigr)\zeta
    +\varepsilon^5\bigl(\Psi M\Psi^\T\dot\Psi+\dot\Psi M^2\bigr)\zeta-\varepsilon^6\Psi M^3\zeta.
\end{align*}

\noindent\emph{(b) The slow part of $\dot V$.} With $H:=I+\varepsilon P_1+\varepsilon^2P_2$ (symmetric),
$\frac{\mathrm{d}}{\mathrm{d}t}[\frac12\zeta^\T H\zeta]=\zeta^\T H\dot\zeta+\frac12\zeta^\T\dot H\zeta$ and
$\dot H=2\varepsilon(M-\bar M)+\varepsilon^2(F-\dot M)$. Insert $\dot\zeta$ and collect by powers of $\varepsilon$.
At order $\varepsilon$: $-\zeta^\T M\zeta+\zeta^\T(M-\bar M)\zeta=-\zeta^\T\bar M\zeta$. At order $\varepsilon^2$:
\begin{align*}
  -\zeta^\T P_1M\zeta+\zeta^\T\Psi^\T\dot\Psi\zeta+\tfrac12\zeta^\T(F-\dot M)\zeta
  &=-\tfrac12\zeta^\T(P_1M+MP_1)\zeta+\tfrac12\zeta^\T\dot M\zeta\\
  &\qquad+\tfrac12\zeta^\T(P_1M+MP_1-\dot M)\zeta=0,
\end{align*}
using $\zeta^\T\Psi^\T\dot\Psi\zeta=\frac12\zeta^\T\dot M\zeta$: the choice of $P_2$ removes the whole second-order
term, and this is where $N_2=0$ enters (otherwise $-\zeta^\T N_2\zeta$ would remain). The remaining terms are
\[
  R_\zeta:=\varepsilon^3\bigl(-\zeta^\T P_2M\zeta+\zeta^\T P_1\Psi^\T\dot\Psi\zeta-\zeta^\T M^2\zeta\bigr)
  +\varepsilon^4\bigl(\zeta^\T P_2\Psi^\T\dot\Psi\zeta-\zeta^\T P_1M^2\zeta\bigr)-\varepsilon^5\zeta^\T P_2M^2\zeta,
\]
plus the coupling $\zeta^\T H\Psi^\T w$, so
$\frac{\mathrm{d}}{\mathrm{d}t}[\frac12\zeta^\T H\zeta]=-\varepsilon\zeta^\T\bar M\zeta+R_\zeta+\zeta^\T H\Psi^\T w$.

\noindent\emph{(c) Per-direction bounds.} For every $v$ and $t$:
\begin{gather*}
  \norm{\Psi v}\le\sqrt\rho\norm{v}_{\bar M},\qquad \norm{\dot\Psi v}\le\sqrt{\rho_D}\norm{v}_{\bar M},\qquad
  \norm{\ddot\Psi v}\le\sqrt{\rho_{DD}}\norm{v}_{\bar M};\\
  \norm{Mv}\le\Psi_M\norm{\Psi v}\le\Psi_M\sqrt\rho\norm{v}_{\bar M};\qquad
  \norm{\bar Mv}\le\lambda_{\max}(\bar M)^{1/2}\norm{v}_{\bar M}\le\Psi_M\norm{v}_{\bar M};\\
  \norm{\Psi^\T\dot\Psi v}\le\Psi_M\sqrt{\rho_D}\norm{v}_{\bar M};\qquad
  \norm{P_1v}\le\tfrac T2\sup_r\norm{(M(r)-\bar M)v}\le\pi_1\norm{v}_{\bar M};\\
  \norm{Fv}\le\norm{P_1Mv}+\norm{MP_1v}\le\tfrac T2\Psi_M^3\sqrt\rho\norm{v}_{\bar M}+\Psi_M^2\pi_1\norm{v}_{\bar M}
  =T\Psi_M^3(\sqrt\rho+\tfrac12)\norm{v}_{\bar M},
\end{gather*}
so $\norm{Sv}\le(T/4)\sup\norm{Fv}=(T^2/8)\Psi_M^3(2\sqrt\rho+1)\norm{v}_{\bar M}$ and
$\norm{P_2v}\le\norm{(M-\bar M)v}+\norm{Sv}\le\pi_2\norm{v}_{\bar M}$; also
\[
  \norm{P_2}\le\Psi_M^2+\tfrac T4\norm{F}_\infty\le\Psi_M^2+\tfrac T4\cdot2\cdot\tfrac{T\Psi_M^2}2\Psi_M^2=p_2.
\]
Hence
$\abs{R_\zeta}\le r_\zeta\norm{\zeta}^2_{\bar M}$,
$\abs{\zeta^\T H\Psi^\T w}=\abs{\ip{\Psi H\zeta}{w}}\le c_1\norm{\zeta}_{\bar M}\norm{w}$, and from (a),
$\beta w^\T\dot w\le\beta\bigl[-(k-d_w)\norm{w}^2+c_2\norm{\zeta}_{\bar M}\norm{w}\bigr]$: each $\zeta$-term of
$\dot w$ is bounded through $\norm{\Psi\zeta}$, $\norm{\dot\Psi\zeta}$, $\norm{\ddot\Psi\zeta}$, $\norm{M\zeta}$ or
$\norm{M^2\zeta}$ (the terms of $c_2$, in the order of (a); the last one is
$\norm{\Psi M^3\zeta}\le\Psi_M\cdot\Psi_M^4\norm{M\zeta}\le\Psi_M^6\sqrt\rho\norm{\zeta}_{\bar M}$), the $w$-terms
through operator norms ($d_w$).

\noindent\emph{(d) Assembly.}
$\dot V\le-\varepsilon\norm{\zeta}^2_{\bar M}+r_\zeta\norm{\zeta}^2_{\bar M}
+(c_1+\beta c_2)\norm{\zeta}_{\bar M}\norm{w}-\beta(k-d_w)\norm{w}^2$. By Young,
$(c_1+\beta c_2)\norm{\zeta}_{\bar M}\norm{w}\le\delta\varepsilon\norm{\zeta}^2_{\bar M}
+(c_1+\beta c_2)^2\norm{w}^2/(4\delta\varepsilon)$, and $(c_1+\beta c_2)^2\le2c_1^2+2\beta^2c_2^2$. With
$\beta=2c_1^2/\delta$ the term $2c_1^2/(4\delta\varepsilon)$ equals $\beta k/4$, and (A) gives
$2\beta^2c_2^2/(4\delta\varepsilon)\le\beta k/4$, so with (B) the coefficient of $\norm{w}^2$ is at most
$-\beta(k-d_w-k/2)\le-\beta k/4$; the coefficient of $\norm{\zeta}^2_{\bar M}$ is
$-\varepsilon(1-\delta-r_\zeta/\varepsilon)=-\varepsilon c_\zeta$. This is the first inequality. Since
$\frac12(1-\varepsilon p_1-\varepsilon^2p_2)\norm{\zeta}^2+\frac\beta2\norm{w}^2\le V
\le\frac12(1+\varepsilon p_1+\varepsilon^2p_2)\norm{\zeta}^2+\frac\beta2\norm{w}^2$ by (D), and
$\zeta^\T\bar M\zeta\ge\lambda_{\min}\norm{\zeta}^2$, $\dot V\le-\gamma V$ follows, and the envelope with $C$ from
$\norm{e}\le\norm{w}+(\varepsilon\Psi_M+\varepsilon^2\Psi_D+\varepsilon^3\Psi_M^3)\norm{\zeta}$ and its converse.
\end{proof}

\begin{rmk}{Remark}{E.0$'$.1}[reading]\mbox{}
\begin{itemize}
  \item \emph{What is sharp and what is not.} The exact rate of (LTV) is $\lambda_{\min}/k\cdot(1+O(k^{-2}))$: the
    first-order averaging of $\dot\zeta=-(1/k)M\zeta$ after the slaving gives $\lambda_{\min}/k$, and the second-order
    correction to the symmetric part of the averaged generator is $N_2=0$ (Table~\ref{tab:thmE:slowregime} measures
    $1-\text{exact}/(\lambda_{\min}/k)\approx29/k^2$ on the ideal; on $\Psi=(\sin t,\cos t)$ the deviation is of
    order $k^{-4}$). Lemma E.0$'$ recovers this rate up to the factor
    $c_\zeta/(1+\varepsilon p_1+\varepsilon^2p_2)\to1-\delta$, and $\delta$ can be taken $\to0$ as $k\to\infty$
    (condition (A) allows $\delta\gtrsim2\sqrt{\rho\rho_{DD}}/k^2$). The overshoot constant $C=O((\rho/\delta)^{1/2})$
    is an artefact of the quadratic form, not of the dynamics.
  \item \emph{Two regimes, two lemmas.} Lemma E.0 needs only $(\alpha,T,\Psi_M,\Psi_D)$ and gives the $\omega^{-2}$ law
    of the fast regime; its rate is quadratic in $\lambda_{\min}$. Lemma E.0$'$ needs $k$ above the threshold
    $k_0\approx(4\rho\rho_{DD}/\delta^2)^{1/4}$, set by the peak factors of $\Psi$ and its two derivatives, and gives
    the sharp rate. On the regressors of Theorem E the derivative peak factors are large---the directions that are
    poorly excited on average are excited in short bursts, so $\rho_{DD}/\rho$ is of the order of the fourth power of
    the effective bandwidth---and $k_0$ is in the hundreds there, while on $\Psi=(\sin t,\cos t)$
    ($\rho=\rho_D=\rho_{DD}=2$) it is about $10$, set by (C) rather than (A). A third-order slaving would replace
    $\rho_{DD}$ by the peak factor of $\dddot\Psi$ at one order higher in $\varepsilon$ and lower $k_0$ further; this is
    not pursued.
  \item \emph{Hypotheses.} Only the periodicity of $\Psi$ is used (through the zero-mean antiderivatives; a
    non-periodic PE regressor would need a windowed $\bar M(t)$ and its derivative). Theorem E's $x_m$ is periodic, so
    $\Psi_G$ is. Theorem E's steps 4--5 go through with this $V$ in place of Lemma E.0's (Corollary E.0$'$.2 below):
    the nonlinear rate becomes $\gamma/4\approx(1-\delta)\lambda_{\min}/(2k)$, while the explicit ball radius carries
    the anisotropy $p'_{\mathrm{lo}}/p'_{\mathrm{hi}}\approx1/\beta$ of the quadratic form and can be smaller than
    Lemma E.0's.
\end{itemize}
\end{rmk}

\paragraph{Numerics} (\code{exp_explicit_rate.py} $\to$ \code{results/explicit_rate/table.md}, table~4; reproduced
in Table~\ref{tab:thmE:slowregime}). The slow-adaptation ideal of Remark E.0a (table~2 of that file) at $g=0.01$
($T=20\pi$): $\lambda_{\min}(\bar M)=2.16\cdot10^{-7}$, $\kappa(\bar M)=7.8\cdot10^3$, $\Psi_M=0.168$,
$\Psi_D=0.957$, $\Psi_{DD}=8.39$, $T\Psi_M^2=1.78$; peak factors $\rho=69.9$, $\rho_D=3.7\cdot10^3$,
$\rho_{DD}=3.7\cdot10^5$ (the poorly excited directions are excited in bursts: $\rho_D/\rho\approx53$ and
$\rho_{DD}/\rho_D\approx99$ are effective squared bandwidths); $\norm{N_2}=3\cdot10^{-21}$. The threshold of (A)--(D)
is $k_0=226$ for $\delta=0.2$ and $321$ for $\delta=0.1$, set by (A): $c_1c_2=2.07$, $0.51$, $0.13$, $0.056$,
$0.032$, $0.014$ at $k=50$, $100$, $200$, $300$, $400$, $600$, the $k^{-2}$ law of $\varepsilon^2\sqrt{\rho\rho_{DD}}$.

\medskip
\noindent\begin{minipage}{\textwidth}
\centering\footnotesize
\captionof{table}[Lemma E.0$'$ in the slow regime]{Lemma E.0$'$ in the slow regime on the slow-adaptation ideal of
Remark E.0a at $g=0.01$ (\code{results/explicit_rate/table.md}, table~4).}
\label{tab:thmE:slowregime}
\begin{tabularx}{\textwidth}{@{}c>{\hsize=0.9\hsize\centering\arraybackslash}X>{\hsize=0.8\hsize\centering\arraybackslash}X>{\hsize=1.1\hsize\centering\arraybackslash}X>{\hsize=0.8\hsize\centering\arraybackslash}X>{\hsize=0.9\hsize\centering\arraybackslash}X>{\hsize=1.5\hsize\centering\arraybackslash}X@{}}
\toprule
$k$ & exact Floquet rate & exact / $(\lambda_{\min}/k)$ & guaranteed rate $\gamma/2$, $\delta=0.2$ ($\delta=0.1$)
  & exact / guaranteed & exact / Lemma E.0 (optimised $\varepsilon$)
  & $\sup\dot V/V$ on the grid, $-\gamma$ ($\delta=0.1$)\\
\midrule
300 & $7.21\cdot10^{-10}$ & 0.9997 & $5.70\cdot10^{-10}$ (below threshold) & 1.26 (---) & $1.33\cdot10^5$ & --- \\
400 & $5.41\cdot10^{-10}$ & 0.9998 & $4.30\cdot10^{-10}$ ($4.84\cdot10^{-10}$) & 1.26 (1.12) & $1.32\cdot10^5$
  & $-1.079\cdot10^{-9}$, $-9.68\cdot10^{-10}$ \\
600 & $3.61\cdot10^{-10}$ & 0.9999 & $2.88\cdot10^{-10}$ ($3.24\cdot10^{-10}$) & 1.25 (1.12) & $1.30\cdot10^5$
  & $-7.20\cdot10^{-10}$, $-6.47\cdot10^{-10}$ \\
\bottomrule
\end{tabularx}
\end{minipage}
\medskip

Above the threshold the guaranteed rate is within 26\% ($\delta=0.2$) and 12\% ($\delta=0.1$) of the exact rate---the
factor $1/(1-\delta)$ of the lemma, with $c_\zeta=0.80$ and $0.90$---while Lemma E.0's optimised rate is
$1.3\cdot10^5$ times smaller; the certificate $\sup\dot V/V\le-\gamma$ holds on the $128$-point grid in every row
where the conditions hold, with $\sup\dot V/V=-1.078\cdot10^{-9}\approx-2\lambda_{\min}/k$ at $k=400$ (attained in
the slow direction at $w=0$). The exact rate deviates from $\lambda_{\min}/k$ by $1.15\%$, $0.29\%$, $0.07\%$ at
$k=50$, $100$, $200$: the $k^{-2}$ law that $N_2=0$ predicts. The overshoot constant is
$C=\sqrt\beta\cdot(1+O(1/k))=27$ ($\delta=0.2$) and $38$ ($\delta=0.1$). Below $k_0$ the lemma makes no claim;
Lemma E.0 does (Remark E.0a; table~2 of that file).

\checkedby{\code{tests/test_excitation.py}: on $\Psi=(\sin t,\cos t)$ the computed $\lambda_{\min}$, $\rho$,
$\rho_D$, $\rho_{DD}$ equal $1/2$, $2$, $2$, $2$, $N_2=0$ to $10^{-12}$, the threshold of (A)--(D) is below $12$ for
$\delta=0.2$, the certificate $\sup\dot V/V\le-\gamma$ holds at $k=20$, $50$, $200$, and the guaranteed rate is below
the exact Floquet rate and within a factor $2$ of it, the ratio decreasing towards $1/(1-\delta)$\newline
(\code{test_lemma_E0prime_sinusoidal_regressor_certificate_and_sharpness}); on the slow-adaptation ideal of the
Theorem E tests at $g=0.01$ and $k=400$, $\delta=0.2$: $\kappa(\bar M)>10^3$, the conditions hold, the certificate
holds on a $128$-point grid, the guaranteed rate is within a factor $1.5$ of the exact rate and exceeds Lemma E.0's
optimised rate by more than $10^4$\newline
(\code{test_lemma_E0prime_pays_for_ill_conditioning_once_on_the_theorem_E_ideal}).}

\begin{thm}{Corollary}{E.0$'$.2}[Theorem E in the slow regime]
Under the hypotheses of Theorem E, suppose the regressor $\Psi_G$ of its step~3 satisfies the conditions (A)--(D) of
Lemma E.0$'$ at $k=k_e$ for some $\delta$, with constants $(\beta,\gamma,p_1,p_2)$ and
$L_M:=\varepsilon\Psi_{G,M}+\varepsilon^2\Psi_{G,D}+\varepsilon^3\Psi_{G,M}^3$ ($\varepsilon=1/k_e$), and that
$\beta L_M^2<1-\varepsilon p_1-\varepsilon^2p_2$. Put
\begin{gather*}
  p'_{\mathrm{lo}}:=\tfrac12\min\{\beta/2,\ 1-\varepsilon p_1-\varepsilon^2p_2-\beta L_M^2\},\\
  p'_{\mathrm{hi}}:=\tfrac12\max\{\beta(1+L_M),\ 1+\varepsilon p_1+\varepsilon^2p_2+\beta L_M(1+L_M)\},\\
  c'_0:=\min\Bigl\{\Bigl[\frac{\gamma p'_{\mathrm{lo}}}{4(p'_{\mathrm{hi}})^{1/2}\tilde C_R}\Bigr]^2,\
  \frac{p'_{\mathrm{lo}}r^2}{\norm{S^{-1}}^2}\Bigr\},\qquad
  r'_0:=\Bigl(\frac{c'_0}{p'_{\mathrm{hi}}}\Bigr)^{1/2},\qquad \lambda':=\frac\gamma4,\qquad
  C':=\Bigl(\frac{2p'_{\mathrm{hi}}}{p'_{\mathrm{lo}}}\Bigr)^{1/2},
\end{gather*}
with $r$, $\tilde C_R$, $S$ as in Theorem E. Then the conclusions of Theorem E hold with $(c_0,\lambda,C)$ replaced by
$((r'_0)^2/2,\lambda',C')$. In the slow regime
$\lambda'=(1-\delta)(1-\varepsilon p_1+O(k_e^{-2}))\lambda_{\min}(\bar M_G)/(2k_e)$,
$\bar M_G:=\langle\Psi_G^\T\Psi_G\rangle$: half the exact rate of the linearised loop, up to the factor $1-\delta$ and
the first-order deficit $\varepsilon p_1$.
\end{thm}

\begin{proof}
Write $V:=\frac12w^\T\mathbf{H}(t)w$ for Lemma E.0$'$'s $V$ in the coordinates $w=(e,\zeta)$ of step~3:
\[
  \mathbf{H}=\begin{bmatrix}\beta I&-\beta L\\-\beta L^\T&H_\zeta+\beta L^\T L\end{bmatrix}
\]
with $e_2=L(t)\zeta$, $\norm{L}\le L_M$ and $H_\zeta:=I+\varepsilon P_1+\varepsilon^2P_2$. Then
$w^\T\mathbf{H}w=\beta\norm{e-L\zeta}^2+\zeta^\T H_\zeta\zeta$, and
$\norm{e-L\zeta}^2\le(1+L_M)(\norm{e}^2+L_M\norm{\zeta}^2)$,
$\norm{e-L\zeta}^2\ge\frac12\norm{e}^2-L_M^2\norm{\zeta}^2$ give
$2p'_{\mathrm{lo}}\norm{w}^2\le w^\T\mathbf{H}w\le2p'_{\mathrm{hi}}\norm{w}^2$. Along the nonlinear loop
$\dot w=[A_0+B(t)]w+\tilde R(t,w)$ with $\norm{\tilde R}\le\tilde C_R\norm{w}^2$ on $\{\norm{w}<r/\norm{S^{-1}}\}$
(step~4), so with Lemma E.0$'$,
\[
  \dot V\le-\gamma V+\ip{\mathbf{H}w}{\tilde R}
  \le-\gamma V+\norm{\mathbf{H}^{1/2}w}\,\norm{\mathbf{H}^{1/2}}\,\tilde C_R\norm{w}^2
  \le-\gamma V+\frac{2(p'_{\mathrm{hi}})^{1/2}\tilde C_R}{p'_{\mathrm{lo}}}V^{3/2},
\]
using $\norm{\mathbf{H}^{1/2}w}=(2V)^{1/2}$, $\norm{\mathbf{H}^{1/2}}=(2p'_{\mathrm{hi}})^{1/2}$ and
$\norm{w}^2\le V/p'_{\mathrm{lo}}$. Hence $\dot V\le-(\gamma/2)V$ wherever
$V^{1/2}\le\gamma p'_{\mathrm{lo}}/(4(p'_{\mathrm{hi}})^{1/2}\tilde C_R)$, and the sublevel set $\{V\le c'_0\}$---inside
the ball where $\tilde R$ is bounded, by the second term of $c'_0$---is forward invariant with
$V(t)\le e^{-(\gamma/2)(t-t_0)}V(t_0)$ on it. It contains $\{\norm{w}\le r'_0\}$ because $V\le p'_{\mathrm{hi}}\norm{w}^2$,
and $\norm{w(t)}\le(p'_{\mathrm{hi}}/p'_{\mathrm{lo}})^{1/2}e^{-(\gamma/4)(t-t_0)}\norm{w(t_0)}$ follows. Step~5 of
Theorem E is unchanged: $\norm{w}^2=2V_\pi$, so $\VOT(t_0)\le(r'_0)^2/2$ puts the state in the ball,
$\norm{e}+W\le\sqrt2\norm{w}$ gives $C'$, and uniqueness and the limits follow as there. The asymptotics of
$\lambda'$ are those of $\gamma$ in Lemma E.0$'$.
\end{proof}

\paragraph{Numerics} (\code{exp_explicit_rate.py} $\to$ \code{results/explicit_rate/table.md}, table~5; reproduced
in Table~\ref{tab:thmE:slowcorollary}). Remark E.1's single-kernel instance ($m=n=1$, $\vartheta^*=(1.5,0.3,0.8)$,
$G=0.3\cdot I$, $x_m=1.5\sin t$) at $k_e=40$: $\lambda_{\min}(\bar M_G)=5.41\cdot10^{-2}$, $\kappa(\bar M_G)=6.7$,
peak factors $\rho=4.8$, $\rho_D=30.3$, $\rho_{DD}=461$, threshold $k_0=24.4$ ($\delta=0.2$) and $33.5$
($\delta=0.1$); $L_J=4.62$ estimated from samples; exact linear rate $1.35\cdot10^{-3}$.

\medskip
\noindent\begin{minipage}{\textwidth}
\centering\small
\captionof{table}[Corollary E.0$'$.2 on the single-kernel instance]{Theorem E in the slow regime on Remark E.1's
single-kernel instance at $k_e=40$ (\code{results/explicit_rate/table.md}, table~5).}
\label{tab:thmE:slowcorollary}
\begin{tabularx}{\textwidth}{@{}l>{\centering\arraybackslash}X>{\centering\arraybackslash}X>{\centering\arraybackslash}X>{\centering\arraybackslash}X@{}}
\toprule
Lyapunov function & explicit ball radius $r_0$ ($w$-coordinates) & nonlinear rate $\lambda$ & overshoot $C$
  & exact linear rate\,/\,$\lambda$\\
\midrule
Lemma E.0, optimised $\varepsilon$ (Theorem E) & $6.15\cdot10^{-6}$ & $2.63\cdot10^{-5}$ & 1.43 & 51 \\
Lemma E.0$'$, $\delta=0.2$ (this corollary)    & $1.53\cdot10^{-6}$ & $4.53\cdot10^{-4}$ & 11.8 & 3.0 \\
Lemma E.0$'$, $\delta=0.1$ (this corollary)    & $8.32\cdot10^{-7}$ & $5.16\cdot10^{-4}$ & 17.1 & 2.6 \\
\bottomrule
\end{tabularx}
\end{minipage}
\medskip

On this well-conditioned regressor ($\kappa\approx7$) Lemma E.0 is only an order of magnitude conservative, and the
corollary trades the certified ball for the rate: $\lambda'$ is within a factor $3$ of the exact linear rate (the
factor $2$ of the Lyapunov argument times $1/(1-\delta)$), while $r'_0$ is four times smaller than Theorem E's $r_0$,
carrying the anisotropy $p'_{\mathrm{lo}}/p'_{\mathrm{hi}}\approx1/\beta$ of the quadratic form. On ill-conditioned
regressors the rate improves by the factor of Remark E.0a (Table~\ref{tab:thmE:slowregime}); the radius comparison
was not computed there.

\checkedby{\code{tests/test_excitation.py}: on the single-kernel instance at $k_e=40$, $\delta=0.2$, from a start
with $\norm{w(0)}\le r'_0$, the implemented closed loop satisfies $\norm{w(t)}\le C'e^{-\lambda't}\norm{w(0)}$ at every
sample over $1.5$ e-folds of the explicit envelope ($\mathrm{d}t=0.005$: the fixed-step loop has a discretisation
floor $\approx6\cdot10^{-6}$ at $\mathrm{d}t=0.02$, above the explicit ball), with $\lambda'$ more than ten times
Theorem E's rate from Lemma E.0 and within a factor $4$ of the exact linear rate\newline
(\code{test_corollary_E0prime_2_envelope_on_the_implemented_loop}).}

\begin{thm}{Corollary}{E.0$'$.3}[the stationary parameter spread of the dithered linearised loop]
Under the hypotheses of Lemma E.0$'$ (conditions (A)--(D) at $(k,\delta)$), consider \textup{(LTV)} with parameter
dither,
\begin{equation}
  \dot e=-ke-\Psi(t)\zeta,\qquad \dd\zeta=\Psi(t)^\T e\dd t+\Sigma\dd B_t,
  \tag{LTV$_\tau$}
\end{equation}
with $\Sigma\in\R^{q\times q'}$ constant and $B$ a standard Brownian motion. Let $V$ be the strict Lyapunov function of
Lemma E.0$'$ and $\mathbf{H}(t)$ its quadratic form ($V=\frac12x^\T\mathbf{H}(t)x$, $x=(e,\zeta)$),
$L(t):=-\varepsilon\Psi+\varepsilon^2\dot\Psi-\varepsilon^3\Psi M$ the slaving map ($e_2=L(t)\zeta$),
$L_M:=\varepsilon\Psi_M+\varepsilon^2\Psi_D+\varepsilon^3\Psi_M^3\ge\sup\norm{L}$, and
\begin{gather*}
  \nu(t):=\tfrac12\tr\bigl(\Sigma^\T\mathbf{H}_{\zeta\zeta}(t)\Sigma\bigr)
  =\tfrac12\tr\bigl(\Sigma^\T(I+\varepsilon P_1+\varepsilon^2P_2+\beta L^\T L)\Sigma\bigr),\\
  \bar\nu:=\sup_t\nu(t)\le\tfrac12\norm{\Sigma}_F^2\bigl(1+\varepsilon p_1+\varepsilon^2p_2+\beta L_M^2\bigr).
\end{gather*}
Then, for every initial law with finite second moment,
\[
  \limsup_{t\to\infty}\E V(t,x_t)\le\frac{\bar\nu}{\gamma},\qquad
  \limsup_{t\to\infty}\E\norm{\zeta_t}^2\le\frac{2\bar\nu}{\gamma(1-\varepsilon p_1-\varepsilon^2p_2)},\qquad
  \limsup_{t\to\infty}\E\norm{w_t}^2\le\frac{2\bar\nu}{\gamma\beta},
\]
with $\gamma$ the rate of Lemma E.0$'$; in the slow regime
$\gamma=2c_\zeta\lambda_{\min}(\bar M)/(k(1+\varepsilon p_1+\varepsilon^2p_2))$, so that
\[
  \limsup\E\norm{\zeta_t}^2\le
  \frac{k\norm{\Sigma}_F^2\,(1+\varepsilon p_1+\varepsilon^2p_2)(1+\varepsilon p_1+\varepsilon^2p_2+\beta L_M^2)}
  {2c_\zeta(1-\varepsilon p_1-\varepsilon^2p_2)\lambda_{\min}(\bar M)}
  =(1+O(\varepsilon))\,\frac{k\norm{\Sigma}_F^2}{2(1-\delta)\lambda_{\min}(\bar M)}.
\]
For comparison, the averaged Langevin dynamics $\dot\zeta=-\frac1k\bar M\zeta+\Sigma\dot B$ has the exact stationary
covariance $\bar C=kX$ with $\bar MX+X\bar M=\Sigma\Sigma^\T$, so $\E\norm{\zeta}^2=k\tr X$, equal to
$\frac k2\tr(\bar M^{-1}\Sigma\Sigma^\T)$ when $\Sigma\Sigma^\T$ commutes with $\bar M$: the bound is within
$(1+O(\varepsilon))/(1-\delta)$ of it when $\bar M$ is isotropic, and pays the condition number of $\bar M$ in the
trace otherwise (it replaces every $1/\lambda_i$ by $1/\lambda_{\min}$), as a scalar Lyapunov function must.
\end{thm}
\status{Proved (It\^o's formula and the pointwise inequality of Lemma E.0$'$). Checked (\code{exp_explicit_rate.py},
table~6; \code{tests/test_excitation.py}) against the exact periodic stationary covariance of \textup{(LTV$_\tau$)},
computed from the periodic Lyapunov equation.}

\begin{proof}
$V(t,x)=\frac12x^\T\mathbf{H}(t)x$ is $C^1$ in $t$ and quadratic in $x$, so by It\^o's formula along
\textup{(LTV$_\tau$)}, with $\tilde\Sigma:=[0;\Sigma]\in\R^{(n+q)\times q'}$,
\begin{multline*}
  \dd V=\bigl[\partial_tV+\nabla_xV\cdot A(t)x\bigr]\dd t+\nabla_xV\cdot\tilde\Sigma\dd B_t
  +\tfrac12\tr\bigl(\tilde\Sigma^\T\mathbf{H}(t)\tilde\Sigma\bigr)\dd t\\
  \le\bigl[-\gamma V+\nu(t)\bigr]\dd t+\nabla_xV\cdot\tilde\Sigma\dd B_t,
\end{multline*}
the inequality being Lemma E.0$'$'s $\dot V\le-\gamma V$ pointwise in $(t,x)$, and
$\tilde\Sigma^\T\mathbf{H}\tilde\Sigma=\Sigma^\T\mathbf{H}_{\zeta\zeta}\Sigma$ with
$\mathbf{H}_{\zeta\zeta}=I+\varepsilon P_1+\varepsilon^2P_2+\beta L^\T L$ from the block form of $\mathbf{H}$
(Corollary E.0$'$.2). The stochastic integral is a martingale because its integrand is linear in $x$ and the linear SDE has
bounded second moments (Gaussian with a bounded covariance), so $\frac{\mathrm{d}}{\mathrm{d}t}\E V\le-\gamma\E V+\bar\nu$ and
$\E V(t)\le e^{-\gamma t}\E V(0)+\bar\nu/\gamma$. The bounds on $\E\norm{\zeta}^2$ and $\E\norm{w}^2$ follow from
$V\ge\frac12(1-\varepsilon p_1-\varepsilon^2p_2)\norm{\zeta}^2$ and $V\ge\frac\beta2\norm{w}^2$; the estimate of
$\bar\nu$ uses $\norm{P_1}\le p_1$, $\norm{P_2}\le p_2$, $\norm{L}\le L_M$. The averaged covariance is the stationary
solution of the Lyapunov equation of the time-invariant averaged system.
\end{proof}

\paragraph{Numerics} (\code{exp_explicit_rate.py} $\to$ \code{results/explicit_rate/table.md}, table~6; reproduced
in Table~\ref{tab:thmE:dither}). $\Sigma=(2\cdot10^{-3})^{1/2}I$; the exact stationary covariance from the periodic
Lyapunov equation (RK4 with $k\cdot\mathrm{d}t\le1/4$; the period mean and the maximum of $\E\norm{\zeta}^2$ coincide
to four digits), the averaged value $k\cdot\tr X$, and the bound:

\medskip
\noindent\begin{minipage}{\textwidth}
\centering\footnotesize
\captionof{table}[Corollary E.0$'$.3: the stationary parameter spread of the dithered (LTV)]{Corollary E.0$'$.3: the
explicit bound on the stationary parameter spread of the dithered (LTV) against the exact periodic stationary
covariance and the averaged value (\code{results/explicit_rate/table.md}, table~6).}
\label{tab:thmE:dither}
\setlength{\tabcolsep}{3pt}
\begin{tabularx}{\textwidth}{@{}>{\hsize=1.1\hsize}Yc>{\hsize=0.75\hsize\centering\arraybackslash}X>{\hsize=1.0\hsize\centering\arraybackslash}X>{\hsize=1.0\hsize\centering\arraybackslash}X>{\hsize=1.45\hsize\centering\arraybackslash}X>{\hsize=0.9\hsize\centering\arraybackslash}X>{\hsize=0.8\hsize\centering\arraybackslash}X@{}}
\toprule
regressor ($\kappa(\bar M)$) & $k$ & $\delta$ & exact $\E\norm{\zeta}^2$ & averaged $k\cdot\tr X$ & bound
  & bound / exact & $1/(1-\delta)$\\
\midrule
Theorem E ideal of table~4 ($7.8\cdot10^3$) & 300 & 0.2 & $1.509\cdot10^6$ & $1.509\cdot10^6$ & $2.64\cdot10^7$
  & 17.5 & 1.25\\
 & 400 & 0.2 / 0.1 & $2.012\cdot10^6$ & $2.012\cdot10^6$ & $3.51\cdot10^7$ / $3.11\cdot10^7$ & 17.4 / 15.5
  & 1.25 / 1.11\\
 & 600 & 0.2 / 0.1 & $3.018\cdot10^6$ & $3.018\cdot10^6$ & $5.23\cdot10^7$ / $4.65\cdot10^7$ & 17.3 / 15.4
  & 1.25 / 1.11\\
$\Psi=(\sin t,\cos t)$ (1) & 20 & 0.2 / 0.1 & $8.025\cdot10^{-2}$ & $8.000\cdot10^{-2}$ & 0.216 / 0.202
  & 2.69 / 2.52 & 1.25 / 1.11\\
 & 50 & 0.2 / 0.1 & $2.001\cdot10^{-1}$ & $2.000\cdot10^{-1}$ & 0.315 / 0.282 & 1.58 / 1.41 & 1.25 / 1.11\\
 & 200 & 0.2 / 0.1 & $8.000\cdot10^{-1}$ & $8.000\cdot10^{-1}$ & 1.051 / 0.935 & 1.31 / 1.17 & 1.25 / 1.11\\
\bottomrule
\end{tabularx}
\end{minipage}
\medskip

The exact stationary spread is the averaged Langevin value up to $3\cdot10^{-4}$ (ideal) and $1.25/k^2$ (sinusoid): on
these instances, where $\Sigma$ is isotropic so that $\Sigma\Sigma^\T$ commutes with $\bar M$, the averaging prediction
of the dithered loop's stationary covariance is exact to $O(k^{-2})$. When $\Sigma\Sigma^\T$ does not commute with
$\bar M$ the deviation is $O(1/k)$ (referee round 7 measured $\pm0.25/k$ on two-dimensional sinusoids with
anisotropic $\Sigma$, the bound still holding with ratios 2.7--5.1). On the isotropic sinusoid the bound exceeds the exact value by a factor that falls toward $1/(1-\delta)$ ($2.7$, $1.6$,
$1.3$ at $k=20$, $50$, $200$ for $\delta=0.2$; the excess is the $O(1/k)$ of the factors $(1\pm\varepsilon p_1)\dots$
and $\beta L_M^2$). On the ideal, where one direction is excited $7.8\cdot10^3$ times less than the others, the
bound is 15--18 times the exact value: it charges every one of the $q=15$ directions the price of the worst,
$(q/\lambda_{\min})/\sum_i(1/\lambda_i)\approx q=15$ (up to $1/(1-\delta)$), the trace-level price of a scalar
Lyapunov function. The exact
$\E\norm{e}^2$ is $\norm{\Sigma}_F^2/(2k)=\tr(\bar MX)/k$ to five digits on both instances (the error rides on its slow
manifold $e\approx\Psi\zeta/k$, and $\tr(\bar MX)=\frac12\norm{\Sigma}_F^2$ by the trace of the Lyapunov equation),
while the bound on $\E\norm{w}^2$ is coarse (factor $1/\beta$): on the sinusoid the exact stationary
$\E\norm{w}^2=\E\norm{e-L\zeta}^2$, computed from the periodic covariance and the slaving map
$L=-\varepsilon\Psi+\varepsilon^2\dot\Psi-\varepsilon^3\Psi M$, is more than three orders of magnitude below the bound
(asserted in the test below).

\checkedby{\code{tests/test_excitation.py}: on $\Psi=(\sin t,\cos t)$ with $\Sigma=(2\cdot10^{-3})^{1/2}I$, the
explicit bound dominates the exact period mean and the exact maximum of $\E\norm{\zeta}^2$ at $k=20$, $50$, $200$
for $\delta=0.1$ and $0.2$, its ratio to the exact value decreases toward $1/(1-\delta)$ (within 6\% at $k=200$), the
grid value of $\bar\nu$ lies below its explicit bound, and the exact value agrees with the averaged $k\cdot\tr X$
within $2/k^2$\newline
{\footnotesize(\code{test_corollary_E0prime_3_dither_covariance_bound_against_the_periodic_lyapunov_solution}).}}

\begin{thm}{Lemma}{E.0$''$}[the slow regime for non-periodic persistently exciting regressors]
Let $\Psi\colon[0,\infty)\to\R^{n\times q}$ be $C^2$ with $\norm{\Psi}\le\Psi_M$, $\norm{\dot\Psi}\le\Psi_D$,
$\norm{\ddot\Psi}\le\Psi_{DD}$, $M:=\Psi^\T\Psi$, and fix a window $T>0$. For a bounded continuous $X$ let
\begin{gather*}
  \langle X\rangle_T(t):=\frac1T\int_0^\infty e^{-s/T}X(t+s)\dd s\\
  \text{(the forward exponential average; $T\tfrac{\mathrm{d}}{\mathrm{d}t}\langle X\rangle_T=\langle X\rangle_T-X$),}
\end{gather*}
and $\bar M:=\langle M\rangle_T$. Assume the exponentially weighted PE condition
\begin{equation}
  \bar M(t)\succeq\lambda I_q\ \text{ for all }t\ge0,\ \text{ with }\lambda>0,
  \tag{PE$_T$}
\end{equation}
which is equivalent to classical PE ($\int_t^{t+T}M\succeq\alpha I$ for all $t$ gives \textup{(PE$_T$)} with
$\lambda\ge e^{-1}\alpha/T$, and \textup{(PE$_T$)} gives PE over windows $T'>T\log(\Lambda/\lambda)$ with level
$T(\lambda-e^{-T'/T}\Lambda)$; proof below). Let $\Lambda:=\sup_t\lambda_{\max}(\bar M(t))\le\Psi_M^2$. Fix symmetric
matrices $\hat M$, $\hat F$ (any; the long-time means of $\bar M$ and $\bar F$ minimise the time-averaged squares of
$\norm{\bar M-\hat M}$, $\norm{\bar F-\hat F}$, not necessarily the suprema below) and define
\begin{gather*}
  P_1:=-2T(\bar M-\hat M)\ \ (\text{so }\dot P_1=2(M-\bar M)),\qquad F:=P_1M+MP_1,\\
  \bar F:=\langle F\rangle_T,\qquad S:=-T(\bar F-\hat F)\ \ (\dot S=F-\bar F),\\
  P_2:=-(M-\bar M)+S,\qquad \bar N_2:=\tfrac12(\bar F-\dot{\bar M})=\tfrac12\bigl(\bar F-(\bar M-M)/T\bigr),
\end{gather*}
and $e_2$, $w$, $V$ exactly as in Lemma E.0$'$
($V=\frac12\zeta^\T(I+\varepsilon P_1+\varepsilon^2P_2)\zeta+\frac\beta2\norm{w}^2$, $\varepsilon=1/k$). The constants
are those of Lemma E.0$'$ with the period mean replaced by the exponential average, all of them computable from the
regressor: the time-local peak factors
\begin{gather*}
  \rho:=\sup_{t,v}\frac{\norm{\Psi(t)v}^2}{v^\T\bar M(t)v},\qquad
  \rho_D:=\sup\frac{\norm{\dot\Psi(t)v}^2}{v^\T\bar M(t)v},\qquad
  \rho_{DD}:=\sup\frac{\norm{\ddot\Psi(t)v}^2}{v^\T\bar M(t)v}\\
  (\text{all }\le\Psi_\cdot^2/\lambda),
\end{gather*}
the operator and per-direction norms
\begin{gather*}
  p_1:=\sup_t\norm{P_1(t)}=2T\sup\norm{\bar M-\hat M},\qquad
  p_2:=\sup_t\norm{P_2(t)}\le\Psi_M^2+T\sup\norm{\bar F-\hat F},\\
  \pi_1:=\sup_{t,v}\frac{\norm{P_1(t)v}}{\norm{v}_{\bar M(t)}}\le2T\bigl(\sqrt\Lambda+\norm{\hat M}/\sqrt\lambda\bigr),\\
  \pi_2:=\sup_{t,v}\frac{\norm{P_2(t)v}}{\norm{v}_{\bar M(t)}}
  \le\Psi_M\sqrt\rho+\sqrt\Lambda+T\bigl(2p_1\Psi_M^2+\norm{\hat F}\bigr)/\sqrt\lambda,\\
  n_2:=\sup_t\norm{\bar N_2(t)}\le\tfrac12\bigl(\sup\norm{F}+2\Psi_M^2/T\bigr)\le p_1\Psi_M^2+\Psi_M^2/T,
\end{gather*}
and, with these $\pi_1$, $\pi_2$, $p_1$, $p_2$, the quantities $c_1$, $c_2$, $d_w$, $r_\zeta$ and
$\beta:=2c_1^2/\delta$ of Lemma E.0$'$ (same formulas), together with
\[
  \delta_2:=\varepsilon n_2/\lambda,\qquad c_\zeta:=1-\delta-\delta_2-r_\zeta/\varepsilon.
\]
Fix $\delta\in(0,1)$. If \textup{(A)} $c_1c_2\le\delta/2$, \textup{(B)} $d_w\le k/4$, \textup{(C)} $c_\zeta>0$ and
\textup{(D)} $\varepsilon p_1+\varepsilon^2p_2<1$, then along \textup{(LTV)}, for all $(e,\zeta,t)$,
\[
  \dot V\le-\varepsilon c_\zeta\,\zeta^\T\bar M(t)\zeta-\frac{\beta k}4\norm{w}^2\le-\gamma V,\qquad
  \gamma:=\min\Bigl\{\frac{2c_\zeta\lambda}{k(1+\varepsilon p_1+\varepsilon^2p_2)},\ \frac k2\Bigr\},
\]
and $\norm{(e(t),\zeta(t))}\le Ce^{-\gamma(t-t_0)/2}\norm{(e(t_0),\zeta(t_0))}$ with the $C$ of Lemma E.0$'$. In the
slow regime the guaranteed rate is
\[
  \gamma/2=(1-\delta-\delta_2)\bigl(1-\varepsilon p_1+O(k^{-2})\bigr)\,\lambda/k,\qquad
  \lambda=\inf_t\lambda_{\min}\bigl(\langle M\rangle_T(t)\bigr):
\]
the exponentially weighted PE level of the window divided by $k$, up to the factor $1-\delta-\delta_2$, where
$\delta_2=O(1/k)$ is the price of the residual second-order term (zero for a periodic $\Psi$ with
$\hat M=\langle M\rangle$ in the limit $T\to\infty$, where Lemma E.0$'$ applies directly).
\end{thm}
\status{Proved (the proof of Lemma E.0$'$ with the changes listed below). Checked (\code{exp_pe_general.py},
\code{tests/test_excitation.py}) on a quasi-periodic regressor and on a regressor driven by a Lorenz orbit, where no
periodic construction applies: the conditions hold from $k=20$ on for the quasi-periodic regressor and from $k=100$ on
for the chaotic one (the derivative peak factors of a chaotic orbit are large), the certificate
$\sup\dot V/V\le-\gamma$ holds on the time grid wherever they hold, and the guaranteed rate is within a factor 1.3--2.5
of the exact decay rate on the quasi-periodic regressor ($k=50$--$200$) and, with the window $T=20$, within a factor
3.5--5 on the chaotic one (shorter windows pay more: 7.5--11 at $T=10$ and 14--138 at $T=5$), where the worst
exponential window is three times less exciting than the long-time average. Referee round 7
(\code{review/REVIEW_round7_E0p3_E0dp.md}) confirmed the proof and the code and found no counterexample on an attack
set of barely exciting, bursty, near-resonant and chaotic regressors.}

\begin{proof}
Follow the proof of Lemma E.0$'$ with $\bar M$, $P_1$, $S$, $P_2$ as defined here. (a) is unchanged (it does not
involve the averages). In (b), $\dot H=2\varepsilon(M-\bar M)+\varepsilon^2(F-\bar F-\dot M+\dot{\bar M})$ exactly,
by $\dot P_1=2(M-\bar M)$ and $\dot P_2=-\dot M+\dot{\bar M}+\dot S=-\dot M+\dot{\bar M}+F-\bar F$; the
order-$\varepsilon$ terms give $-\zeta^\T\bar M\zeta$ as before, and the order-$\varepsilon^2$ terms give
$-\zeta^\T P_1M\zeta+\frac12\zeta^\T\dot M\zeta+\frac12\zeta^\T(F-\bar F-\dot M+\dot{\bar M})\zeta
=\frac12\zeta^\T(\dot{\bar M}-\bar F)\zeta=-\zeta^\T\bar N_2\zeta$, since $\zeta^\T P_1M\zeta=\frac12\zeta^\T F\zeta$:
the periodic cancellation $N_2=0$ is replaced by the residual $-\varepsilon^2\zeta^\T\bar N_2\zeta$, and $R_\zeta$ has
the same form, so
$\frac{\mathrm{d}}{\mathrm{d}t}[\frac12\zeta^\T H\zeta]=-\varepsilon\zeta^\T(\bar M+\varepsilon\bar N_2)\zeta+R_\zeta+\zeta^\T H\Psi^\T w$.
In (c), every per-direction bound is taken in the time-local norm $\norm{v}_{\bar M(t)}$:
$\norm{\Psi(t)v}\le\sqrt\rho\norm{v}_{\bar M(t)}$ and likewise for $\dot\Psi$, $\ddot\Psi$ by the definition of the
peak factors; $\norm{M(t)v}\le\Psi_M\sqrt\rho\norm{v}_{\bar M(t)}$;
$\norm{\bar M(t)v}\le\sqrt\Lambda\norm{v}_{\bar M(t)}$;
$\norm{P_1(t)v}=2T\norm{(\bar M(t)-\hat M)v}\le2T(\sqrt\Lambda+\norm{\hat M}/\sqrt\lambda)\norm{v}_{\bar M(t)}$, using
$\sqrt\lambda\norm{v}\le\norm{v}_{\bar M(t)}$ (this is where the exponential average pays off: $P_1(t)$ is a function
of $\bar M(t)$ alone, with no integral over other times); for $\bar F(t)v$ and $S(t)v$, which do involve other times
$r\ge t$,
$\norm{F(r)v}\le\norm{P_1(r)}\norm{M(r)v}+\norm{M(r)}\norm{P_1(r)v}\le2p_1\Psi_M^2\norm{v}\le(2p_1\Psi_M^2/\sqrt\lambda)\norm{v}_{\bar M(t)}$,
an average of such bounds bounds $\norm{\bar F(t)v}$, and $\norm{S(t)v}\le T(\norm{\bar F(t)v}+\norm{\hat Fv})$,
whence the bound on $\pi_2$; the stated bounds on $p_2$ and $n_2$ are $\norm{P_2}\le\norm{M-\bar M}+\norm{S}$ and
$\norm{\bar N_2}\le\frac12(\norm{\bar F}+\norm{\dot{\bar M}})$ with $\norm{\dot{\bar M}}=\norm{\bar M-M}/T\le2\Psi_M^2/T$
and $\norm{\bar F}\le\sup\norm{F}\le2p_1\Psi_M^2$. The bounds $\abs{R_\zeta}\le r_\zeta\norm{\zeta}^2_{\bar M}$,
$\abs{\zeta^\T H\Psi^\T w}\le c_1\norm{\zeta}_{\bar M}\norm{w}$ and
$\beta w^\T\dot w\le\beta[-(k-d_w)\norm{w}^2+c_2\norm{\zeta}_{\bar M}\norm{w}]$ follow with the same formulas (the
suprema over $v$ at the current time $t$ of the norms that appear). In (d), the extra term is bounded by
$\varepsilon^2\abs{\zeta^\T\bar N_2\zeta}\le\varepsilon^2n_2\norm{\zeta}^2\le(\varepsilon^2n_2/\lambda)\zeta^\T\bar M\zeta=\varepsilon\delta_2\zeta^\T\bar M\zeta$,
which turns $1-\delta-r_\zeta/\varepsilon$ into $c_\zeta=1-\delta-\delta_2-r_\zeta/\varepsilon$; the rest of the
assembly, the sandwich
$\frac12(1-\varepsilon p_1-\varepsilon^2p_2)\norm{\zeta}^2+\frac\beta2\norm{w}^2\le V\le\frac12(1+\varepsilon p_1+\varepsilon^2p_2)\norm{\zeta}^2+\frac\beta2\norm{w}^2$
and $\zeta^\T\bar M(t)\zeta\ge\lambda\norm{\zeta}^2$ give $\dot V\le-\gamma V$. The equivalence of \textup{(PE$_T$)}
with classical PE: if $\int_t^{t+T}M\succeq\alpha I$ then
$\bar M(t)\succeq\frac1T\int_0^Te^{-s/T}M(t+s)\dd s\succeq\frac{e^{-1}}T\int_t^{t+T}M\succeq(e^{-1}\alpha/T)I$;
conversely \textup{(PE$_T$)} gives, for every $T'>0$,
\[
  \int_t^{t+T'}M\succeq\int_t^{t+T'}Me^{-(r-t)/T}\dd r\succeq T\bigl(\bar M(t)-e^{-T'/T}\bar M(t+T')\bigr)
  \succeq T(\lambda-e^{-T'/T}\Lambda)I,
\]
positive for $T'>T\log(\Lambda/\lambda)$ (the identity
$\int_t^{t+T'}e^{-(r-t)/T}M(r)\dd r=T[\bar M(t)-e^{-T'/T}\bar M(t+T')]$ follows from the definition of the exponential
average).
\end{proof}

\begin{rmk}{Remark}{E.0$''$.1}[reading]\mbox{}
\begin{itemize}
  \item \emph{What changes and what does not.} The construction, the conditions and the rate formula are those of
    Lemma E.0$'$; the period mean becomes the exponential average over a window, the peak factors become time-local,
    and one new term appears: the residual $\bar N_2$, whose price $\delta_2=\varepsilon n_2/\lambda$ vanishes like
    $1/k$. For a periodic regressor the exponential average with a finite window keeps residual oscillations of
    relative size $O(\text{period}/T)$, so Lemma E.0$'$ is sharper there; Lemma E.0$''$ is the version for references
    that are not periodic --- quasi-periodic (several incommensurate frequencies) or chaotic (a reference along a
    strange attractor, as in the Lorenz-tracking experiments of Part~I), where only PE is available.
  \item \emph{The window is a design parameter.} The lemma holds for every $T>0$; the guaranteed rate is
    $(1-\delta-\delta_2)\lambda_T/k$ with $\lambda_T=\inf\lambda_{\min}(\langle M\rangle_T)$, which increases toward
    the long-time excitation level as $T$ grows, against $p_1\propto T$, $n_2\propto T$ and the threshold in $k$,
    which grow with $T$. The numerics show the optimum at a window of a few characteristic times of the regressor.
  \item \emph{Exact rate.} For a non-periodic regressor the exact rate is the top Lyapunov exponent of
    \textup{(LTV)}, which the certificate does not need: the lemma is a pointwise inequality in $(t,e,\zeta)$,
    checked on a grid in $t$ by the same generalised eigenvalue problem as Lemma E.0$'$.
\end{itemize}
\end{rmk}

\paragraph{Numerics} (\code{exp_pe_general.py} $\to$ \code{results/explicit_rate/general_pe.md}). Two regressors
that are not periodic, on $[0,300]$ (grid step $0.005$; the exponential averages formed from an extended grid; the
certificate $\sup\dot V/V$ over 400 times and all states; the exact rate the top Lyapunov exponent of \textup{(LTV)}
over 200 time units), with the measured constants (the grid suprema of $\norm{P_1}$, $\norm{P_2}$, the per-direction
norms and $\norm{\bar N_2}$, computable constants of the regressor like the peak factors).

\smallskip
\noindent\emph{Quasi-periodic $\Psi=(\sin t,\cos\sqrt2\,t)$} ($q=2$; the long-time mean of $M$ is $\frac12I$, and
the exact decay rate is $\lambda_\infty/k=0.5/k$ to within 1\%, 0.6\% at $k=20$ and 0.8\% at $k=400$): the window
average has $\lambda_T=0.297$, $0.391$, $0.444$ for $T=2\pi$, $4\pi$, $8\pi$ ($\rho=4.0$, $\rho_D\approx6.1$--$6.4$,
$\rho_{DD}\approx10.0$, $p_1\approx2.9$, $p_2=8.0$--$9.2$, $\pi_1=4.3$--$4.7$, $\pi_2=11.6$--$13.0$,
$\norm{\bar N_2}=0.71$, $0.35$, $0.17$); the conditions hold from $k=20$ on for both $\delta$, the certificate holds
at every row, and (Table~\ref{tab:thmE:generalpe})

\medskip
\noindent\begin{minipage}{\textwidth}
\centering\footnotesize
\captionof{table}[Lemma E.0$''$ on the quasi-periodic regressor]{Lemma E.0$''$ on the quasi-periodic regressor
$\Psi=(\sin t,\cos\sqrt2\,t)$: the guaranteed rate against the exact decay rate for two windows
(\code{results/explicit_rate/general_pe.md}).}
\label{tab:thmE:generalpe}
\begin{tabular}{@{}ccccccc@{}}
\toprule
$T$ & $k$ & $\delta_2$ & guaranteed $\gamma/2$ ($\delta=0.2$ / $0.1$) & exact & exact/guaranteed & $\lambda_T/k$\\
\midrule
$4\pi$ & 20 & 0.045 & 0.0100 / 0.0117 & 0.0249 & 2.49 / 2.13 & 0.0195\\
$4\pi$ & 50 & 0.018 & 0.00557 / 0.00631 & 0.00995 & 1.78 / 1.58 & 0.00781\\
$4\pi$ & 200 & 0.004 & 0.00153 / 0.00172 & 0.00249 & 1.63 / 1.45 & 0.00195\\
$8\pi$ & 20 & 0.020 & 0.0118 / 0.0137 & 0.0249 & 2.10 / 1.81 & 0.0222\\
$8\pi$ & 200 & 0.002 & 0.00174 / 0.00196 & 0.00249 & 1.43 / 1.27 & 0.00222\\
\bottomrule
\end{tabular}
\end{minipage}
\medskip

\noindent($T=2\pi$: $3.69\to2.16$ for $\delta=0.2$ and $3.10\to1.92$ for $\delta=0.1$ from $k=20$ to $200$.) The
ratio exact/guaranteed is the product of the window loss $\lambda_\infty/\lambda_T$ ($1.69$, $1.28$, $1.13$ for the
three windows), the factor $1/(1-\delta-\delta_2)$ and the $O(1/k)$ factors; the explicit closed-form constants
satisfy the conditions from $k=50$ on and give rates below the measured-constant ones by 22--79\% at $k=50$ (the
longer the window, the larger the deficit), 5--19\% at $k=100$, at most 5\% at $k=200$ and about 1\% at $k=400$. All
constants are suprema over $[0,300]$ (the exponential averages are formed on an extended grid): on the quasi-periodic
regressor they are identical to four digits on $[0,10^4]$ (closed form); on the Lorenz regressor below, the suprema
over $[0,1500]$ are larger by up to 14\% ($p_1$ $5.31\to5.67$, $p_2$ $19.4\to20.7$, $\pi_2$ $49.4\to51.9$,
$\norm{\bar N_2}$ $0.42\to0.48$ at $T=20$), which lowers the guaranteed rate at $k=100$ by 3\% and keeps the
certificate (referee round 7): the constants of a chaotic regressor are those of the stretch of the attractor that
the run visits.

\smallskip
\noindent\emph{Chaotic $\Psi=(x_1/20,(x_3-25)/20)$ along a Lorenz orbit} ($\sigma=10$, $r=28$, $b=8/3$, traversed
five times slower than standard Lorenz time, $\dot\Psi$ by the chain rule through the vector field; $q=2$): the exact
decay rates are $0.00759$, $0.00304$, $0.00152$, $0.00076$, $0.00038$ at $k=20$, $50$, $100$, $200$, $400$ (an
effective excitation level $0.15$), while the worst exponential window has $\lambda_T=0.015$, $0.024$, $0.051$ for
$T=5$, $10$, $20$ ($\Lambda=0.55$, $0.36$, $0.24$; $\rho\approx20$--$26$, $\rho_D\approx110$--$160$,
$\rho_{DD}\approx2000$--$3100$; $\norm{\bar N_2}=1.05$, $0.34$, $0.42$: the residual does not follow a monotone rule
in $T$). The conditions hold from $k=100$ on (condition (A) fails below: the derivative peak factors of a chaotic
orbit are large), the certificate holds at every admissible row, and at $T=20$ the guaranteed rate is $0.00030$ /
$0.00035$ ($k=100$), $0.00018$ / $0.00021$ ($k=200$) against the exact $0.00152$ and $0.00076$:
exact/guaranteed $=5.0$ / $4.3$ and $4.2$ / $3.7$, of which the window loss $\lambda_\infty/\lambda_T\approx3$ is the
dominant factor. For a chaotic reference the uniform-in-time construction pays the least exciting stretch of the
attractor; a longer window ($T=20$ against the slowed Lorenz loop time $\approx3.75$) recovers part of it
($\lambda_T=0.015\to0.051$ from $T=5$ to $20$) at the price of a larger $p_2$ and $\pi_2$ ($9.9\to19.4$ and
$13.7\to49.4$). An exploratory regressor with components $\tanh(x_i/8)$ traversed at standard Lorenz speed (not in
the table) had $\rho_{DD}$ of order $10^6$--$10^7$ and a threshold beyond $k=400$: the slow regime requires the
reference to be slow relative to the loop, and the time scale of a chaotic reference is the designer's choice.

\checkedby{\code{tests/test_excitation.py}: on $\Psi=(\sin t,\cos\sqrt2\,t)$ with $T=4\pi$, the differential
identities $T\dot{\bar M}=\bar M-M$ and $\dot P_1=2(M-\bar M)$ on the grid ($10^{-4}$), the conditions at $k=20$,
$50$, $200$ for $\delta=0.1$, $0.2$ with $\delta_2\le5\%$, the certificate $\sup\dot V/V\le-\gamma$ on 400 times, and
the guaranteed rate below the exact decay rate (top Lyapunov exponent over 200 time units) within a factor 2.6, the
ratio improving with $k$ (\code{test_lemma_E0dprime_quasi_periodic_regressor_certificate_and_rate}).}

\begin{thm}{Theorem}{E}[with explicit constants]
Let $\Psi_G:=m^{-1/2}\Psi(I_m\otimes G^{1/2})\in\R^{n\times mp}$, and let $(\alpha_G,T,\Psi_{G,M},\Psi_{G,D})$ be PE
data for it: $\int_t^{t+T}\Psi_G^\T\Psi_G\dd r\succeq\alpha_GI$, $\norm{\Psi_G}\le\Psi_{G,M}$,
$\norm{\dot\Psi_G}\le\Psi_{G,D}$. One may take $\alpha_G:=\alpha\lambda_{\min}(G)/m$,
$\Psi_{G,M}:=m^{-1/2}\lambda_G^{1/2}\sup\norm{\Psi}$ and $\Psi_{G,D}:=m^{-1/2}\lambda_G^{1/2}\sup\norm{\dot\Psi}$;
$\dot\Psi$ is bounded because $x_m\in C^1$ with bounded derivative and the Gaussian's mixed derivatives are bounded
for $s\ge s_{\min}$. Let $(\varepsilon,\gamma,p_{\mathrm{lo}},p_{\mathrm{hi}},g_V)$ be the constants of Lemma E.0 for
the data $(k_e,\alpha_G,T,\Psi_{G,M},\Psi_{G,D})$, with any admissible $\varepsilon$ (for example the closed form of
E.0(ii)), where $g_V:=1+\varepsilon(2Q_{G,M}+T\Psi_{G,M})$ bounds $\norm{\partial V/\partial w}/\norm{w}$ (the
Hessian of $V$ is $I$ plus a matrix of norm at most $\varepsilon T\Psi_{G,M}+2\varepsilon Q_{G,M}$) and
$Q_{G,M}:=T^2\Psi_{G,M}^2/2$. Let $r:=\min(\rho,\delta)$, let $L_J$ be a Lipschitz constant, uniform in $t$, of the Jacobian of the closed-loop field
$F$ of step~2 on the ball $D$, and let $S:=\operatorname{diag}\bigl(I_n,\,m^{-1/2}(I_m\otimes G^{-1/2})\bigr)$ be the
similarity to the coordinates $w:=(e,\zeta)$ of step~3. Put
\begin{gather*}
  \tilde C_R:=\norm{S}\,\frac{L_J}{2}\,\norm{S^{-1}}^2,\qquad
  r_w:=\min\Bigl\{\frac{r}{\norm{S^{-1}}},\ \frac{\gamma}{2g_V\tilde C_R}\Bigr\},\qquad
  r_0:=r_w\Bigl(\frac{p_{\mathrm{lo}}}{p_{\mathrm{hi}}}\Bigr)^{1/2},\\
  \lambda:=\frac{\gamma}{4p_{\mathrm{hi}}},\qquad c_0:=\frac{r_0^2}{2},\qquad
  C:=\Bigl(\frac{2p_{\mathrm{hi}}}{p_{\mathrm{lo}}}\Bigr)^{1/2}\quad(\le\sqrt6\text{ for the closed-form }\varepsilon).
\end{gather*}
If $\VOT(t_0)\le c_0$, then:
\begin{itemize}
  \item the optimal assignment $\pi$ at $t_0$ is unique, and it is the unique optimal assignment for all large $t$;
  \item $\norm{e(t)}+W(\mu_t,\mu^*)\le Ce^{-\lambda(t-t_0)}\bigl(\norm{e(t_0)}+W(\mu_{t_0},\mu^*)\bigr)$ for
    $t\ge t_0$;
  \item $\vartheta_i(t)\to\vartheta^*_{\pi(i)}$. The labelled error
    \[
      E_{\mathrm{lab}}(t):=\Bigl(\frac1m\sum_i\norm{\vartheta_i(t)-\vartheta^*_i}^2_{G^{-1}}\Bigr)^{1/2}
      \quad\text{tends to}\quad
      \Bigl(\frac1m\sum_i\norm{\vartheta^*_{\pi(i)}-\vartheta^*_i}^2_{G^{-1}}\Bigr)^{1/2},
    \]
    which is non-zero for $\pi\neq\mathrm{id}$.
\end{itemize}
Every constant is an explicit function of $(k_e,\alpha_G,T,\Psi_{G,M},\Psi_{G,D},m,G,r,L_J)$, hence uniform over the
class of regressors with given PE data $(T,\alpha,\sup\norm{\Psi},\sup\norm{\dot\Psi})$. With the closed-form
$\varepsilon$, $\lambda\ge\frac13\cdot\min\{k_e/4,\,\varepsilon\alpha_G/2\}$. The constants are conservative, in a
quantified way (Remark E.0a); their role is to make the statement self-contained. $L_J$ is defined as a supremum; a
closed-form bound follows from the third-derivative suprema of the Gaussian and is not computed here
(\code{excitation.jacobian_lipschitz_estimate} estimates it from samples for the tests). The previous version proved
this theorem with existence constants, conditional on Morgan--Narendra~\cite[Thm.~1]{S8}; Lemma E.0 removes that
citation.
\end{thm}
\status{Proved, self-contained (Lemma E.0 replaces the citation of Morgan--Narendra~\cite{S8}); linearisation
checked; the strict Lyapunov inequalities checked pointwise (E.0$'$ by a certificate over all states) and the explicit
envelope checked on the linearised and on the implemented loop; convergence observed. The rate of E.0 is conservative
in a quantified way (Remark E.0a); the rate of E.0$'$ is within 12\% of the exact rate above its threshold. Novelty
(literature audit): partially anticipated~\cite{R49}, \cite{R54}; the strict Lyapunov function is of the
Mazenc--Malisoff type.}

\begin{proof}\mbox{}
\begin{enumerate}[label=\arabic*.]
  \item \emph{Symmetry.} Relabelling the network, $(e,\vartheta)\mapsto(e,\vartheta\circ\sigma)$, maps solutions to
    solutions. It fixes $e$, is an isometry of the matched metric, and conjugates the dynamics near
    $(0,\vartheta^*\circ\pi)$ to those near $(0,\vartheta^*)$. So it suffices to treat $\pi=\mathrm{id}$.

  \item \emph{Domain and linearisation.} Let $r:=\min(\rho,\delta)$ and
    $D:=\{z=(e,\tilde\vartheta)\in\R^{n+mp}:\norm{z}_2<r\}$.
    \begin{itemize}
      \item For $z\in D$ we have $\norm{e}<\rho$, so $f=\fhat^*$.
      \item Also $\norm{\tilde\vartheta_i}<\delta$ for every $i$, so $\vartheta_i\in\operatorname{int}\ThKe$ and
        $\Proj=\mathrm{id}$.
    \end{itemize}
    On $[0,\infty)\times D$ the closed loop is
    \[
      F(t,z)=\Bigl(-k_ee+\fhat^*(x)-\fhat(x),\ \bigl(G\nabla_\vartheta\Phi(x;\vartheta_i)^\T e\bigr)_i\Bigr),
      \qquad x=x_m(t)+e.
    \]
    \begin{itemize}
      \item It is $C^1$ in $(t,z)$, since $x_m\in C^1$, and $C^\infty$ in $z$.
      \item Its derivatives through order two are bounded uniformly in $t$, because the Gaussian's derivatives through
        order three are bounded uniformly in $x$ for $s\ge s_{\min}$.
      \item Its Jacobian at $0$ is
        \[
          A(t)=\begin{pmatrix}-k_eI&-\frac1m\Psi\\[2pt](I_m\otimes G)\Psi^\T&0\end{pmatrix},
        \]
        and the remainder is $O(\norm{z}^2)$.
      \item The Jacobian and the identity $QA+A^\T Q=\operatorname{diag}(-k_eI,0)$, with
        $Q=\operatorname{diag}\bigl(I/2,\,(I\otimes G^{-1})/(2m)\bigr)$, are checked numerically by\newline
        \code{test_theorem_E_linearisation_and_lyapunov_identity}.
    \end{itemize}

  \item \emph{LTV stability, with explicit constants.} Put $w:=(e,\zeta)$ with
    $\zeta:=m^{-1/2}(I_m\otimes G^{-1/2})\tilde\vartheta$, that is $w=Sz$, and $\Psi_G:=m^{-1/2}\Psi(I\otimes G^{1/2})$.
    \begin{itemize}
      \item The system becomes $\dot e=-k_ee-\Psi_G\zeta$ and $\dot\zeta=\Psi_G^\T e$: this is (LTV) of Lemma E.0
        with $k=k_e$ and $\Psi=\Psi_G$, that is, $\dot w=[A_0+B(t)]w$ with $A_0=\operatorname{diag}(-k_eI,0)$ and $B$
        skew-symmetric, bounded, with bounded derivative.
      \item The excitation transfers: $\int\Psi_G^\T\Psi_G\dd r\succeq(\alpha\lambda_{\min}(G)/m)I=\alpha_GI$, and
        $\norm{\Psi_G}\le\Psi_{G,M}$, $\norm{\dot\Psi_G}\le\Psi_{G,D}$.
      \item Lemma E.0 gives the strict Lyapunov function $V(t,w)$ with
        $p_{\mathrm{lo}}\norm{w}^2\le V\le p_{\mathrm{hi}}\norm{w}^2$ and $\dot V\le-\gamma\norm{w}^2$ along the
        linear system, hence
        \[
          \norm{w(t)}\le\Bigl(\frac{p_{\mathrm{hi}}}{p_{\mathrm{lo}}}\Bigr)^{1/2}
          e^{-(\gamma/(2p_{\mathrm{hi}}))(t-t_0)}\norm{w(t_0)}:
        \]
        uniform exponential stability with explicit overshoot and rate.
      \item The $w$-coordinates are the coordinates of the matched Lyapunov function:
        $\norm{\zeta}^2=\frac1m\sum_i\norm{\tilde\vartheta_i}^2_{G^{-1}}=W_\pi^2$, so
        $\norm{w}^2=\norm{e}^2+W_\pi^2=2V_\pi$.
      \item The bound on $\dot\Psi$ cannot be dropped. $\Psi=\sin t^2$ is bounded and PE, but for
        $\dot e=-e-\Psi\zeta$, $\dot\zeta=\Psi e$, from $w(t_0)=(e,\zeta)=(0,1)$, a numerical computation gives the
        ratio $\norm{w(t_0+200)}/\norm{w(t_0)}$ as $0.9715$, $0.9942$ and $0.9980$ for $t_0=5$, $20$ and $50$, and
        part~(f) of the proof of Lemma E.0 shows that the decay over a fixed horizon tends to $0$ as $t_0\to\infty$.
        So that system is not UES. Consistently, the rate of Lemma E.0 tends to $0$ as $\Psi_D\to\infty$.
    \end{itemize}

  \item \emph{Nonlinear system, with an explicit ball.} On the closure of $D$ (to which the Lipschitz bound extends)
    the field is $F(t,z)=A(t)z+R(t,z)$ with $\norm{R(t,z)}\le(L_J/2)\norm{z}^2$ (Taylor with the Lipschitz
    Jacobian). In $w$-coordinates,
    $\dot w=[A_0+B(t)]w+\tilde R(t,w)$ with $\tilde R(t,w):=S\,R(t,S^{-1}w)$, so
    $\norm{\tilde R(t,w)}\le\tilde C_R\norm{w}^2$ on $\{\norm{w}<r/\norm{S^{-1}}\}\subset S(D)$. Along the nonlinear
    system, since $\norm{\partial V/\partial w}\le g_V\norm{w}$,
    \[
      \dot V\le-\gamma\norm{w}^2+g_V\tilde C_R\norm{w}^3\le-\frac{\gamma}{2}\norm{w}^2
      \qquad\text{whenever }\norm{w}\le r_w.
    \]
    The sublevel set $\{V\le p_{\mathrm{lo}}r_w^2\}$ is contained in $\{\norm{w}\le r_w\}$, so it is forward
    invariant, and it contains the ball $\{\norm{w}\le r_0\}$, since $V\le p_{\mathrm{hi}}\norm{w}^2$. On it
    $\dot V\le-(\gamma/(2p_{\mathrm{hi}}))V$, so for $\norm{w(t_0)}\le r_0$,
    \[
      \norm{w(t)}\le\Bigl(\frac{p_{\mathrm{hi}}}{p_{\mathrm{lo}}}\Bigr)^{1/2}e^{-\lambda(t-t_0)}\norm{w(t_0)},
      \qquad\lambda=\frac{\gamma}{4p_{\mathrm{hi}}}.
    \]
    These trajectories stay in $D$, where the closed loop equals $F$. (Khalil~\cite[Thm.~4.13]{khalil2002} gives the
    same conclusion with existence constants; the explicit form is what Lemma E.0 adds.)

  \item \emph{Transfer to $(\norm{e},W)$.} Let $\pi$ be optimal at $t_0$. Then
    $\norm{w(t_0)}^2=2V_\pi(t_0)=2\VOT(t_0)\le2c_0=r_0^2$, so step~4 applies. Moreover
    $\norm{e}+W\le\norm{e}+W_\pi=\norm{e}+\norm{\zeta}\le\sqrt2\norm{w}$ at every $t$, and
    $\norm{w(t_0)}\le\norm{e(t_0)}+\norm{\zeta(t_0)}=\norm{e(t_0)}+W(t_0)$. This gives the estimate with
    $C=(2p_{\mathrm{hi}}/p_{\mathrm{lo}})^{1/2}$.

    \emph{Uniqueness.} Suppose $V_\pi,V_{\pi'}\le c_0$ with $\pi\neq\pi'$. Then the state lies in the invariant balls
    of step~4 around both $(0,\vartheta^*\circ\pi)$ and $(0,\vartheta^*\circ\pi')$. The solution would converge to
    two distinct equilibria, which is impossible.

    \emph{Limits.} Hence $z(t)\to0$, that is, $\vartheta_i(t)\to\vartheta^*_{\pi(i)}$. For $\pi'\neq\pi$,
    \[
      V_{\pi'}(t)\to\frac1{2m}\sum_i\norm{\vartheta^*_{\pi(i)}-\vartheta^*_{\pi'(i)}}^2_{G^{-1}}\ge\frac{d^2}2>0,
    \]
    while $V_\pi(t)\to0$. So $\pi$ is the unique optimal assignment for all large $t$, and $E_{\mathrm{lab}}$
    converges as stated.\qedhere
\end{enumerate}
\end{proof}
\checkedby{\code{tests/test_transport.py}, \code{tests/test_excitation.py} and\newline
\code{experiments/exp_pe_convergence.py}.
\begin{itemize}[nosep]
  \item Linearisation: the Jacobian of the implemented loop equals $A(t)$, and
    $QA+A^\T Q=\operatorname{diag}(-k_eI,0)$\newline
    (\code{test_theorem_E_linearisation_and_lyapunov_identity}).
  \item Explicit constants: on the implemented loop, from a start inside the explicit ball,
    $\norm{w(t)}\le(p_{\mathrm{hi}}/p_{\mathrm{lo}})^{1/2}e^{-\lambda t}\norm{w(0)}$ at every sample ($L_J$ estimated
    from samples)\newline
    (\code{test_theorem_E_explicit_nonlinear_envelope_on_the_implemented_loop}).
  \item Convergence: the matched error decays and the labelled error does not (qualitative); sliding-window
    excitation and fitted rates\newline
    (\code{test_theorem_E_matched_error_converges_labelled_does_not}, \code{exp_pe_convergence.py}).
  \item Not checked in the repository: the sharpness of the explicit rate beyond the tables of Remark E.0a;
    Corollary E.2 and Remark E.3.
\end{itemize}}

\begin{rmk}{Remark}{E.1}[the rate]\mbox{}
\begin{itemize}
  \item \textbf{What converges.} ``The parameters converge'' is false for the labelled error and true for the matched
    one.
  \item \textbf{The rate is not set by the PE level.} $\lambda>0$ exists for each admissible $\Psi$, but it depends
    on $\Psi(\cdot)$, not only on its PE level. No positive function of $\alpha\lambda_{\min}(G)/m$, $k_e$, $T$ and
    $\sup\norm{\Psi_G}$ alone bounds it below.
    \begin{itemize}
      \item For $m=n=1$, $\vartheta^*=(1.5,0.3,0.8)$, $G=I$, $k_e=3$ and $x_m=1.5\sin\omega t$ ($T=2\pi$), these four
        quantities do not depend on $\omega$. The exact Floquet rate, computed numerically, falls from
        $4.7\times10^{-2}$ ($\omega=1$) to $2.9\times10^{-5}$ ($\omega=64$), like $\omega^{-2}$.
      \item A lower bound must also involve $\sup\norm{\dot\Psi_G}$ (step~3). Lemma E.0 gives one; it reproduces the
        $\omega^{-2}$ law (Remark E.0a).
      \item Exact Floquet rates, computed numerically, are also not monotone in $G$ or in $k_e$, and they decrease for
        large $G$. When $k_e$ is large relative to the reference bandwidth, they approach
        $\lambda_{\min}\langle\Psi_G^\T\Psi_G\rangle/k_e$.
    \end{itemize}
  \item \textbf{The numbers} (\code{results/pe_convergence}; $m=3$, $n=2$, six runs):
    \begin{center}
      \small
      \begin{tabularx}{0.9\linewidth}{Yl}
        \toprule
        quantity & range\\
        \midrule
        min over sliding $60$\,s windows of $\lambda_{\min}\bigl(\int_t^{t+60}\Psi^\T\Psi\dd r\bigr)$
          & $5.9\times10^{-4}$ to $2.7\times10^{-2}$\\
        fitted decay rate of the matched error, second half of the run
          & $2.3\times10^{-3}$ to $4.3\times10^{-3}$\,s$^{-1}$\\
        rank correlation between the two & $0.60$\\
        \bottomrule
      \end{tabularx}
    \end{center}
\end{itemize}
The convergence test is therefore qualitative. The machine-checked part of the proof is the linearisation.
\end{rmk}

\begin{thm}{Corollary}{E.2}[uniform-in-time exponential merging under PE]
Take two solutions of the same closed loop with $\VOT(0)\le c_0$ and $V'_{\mathrm{OT}}(0)\le c_0$. Then
\begin{equation}\label{eq:thmE:merging}
  \norm{x_t-x'_t}+W(\mu_t,\mu'_t)\le Ce^{-\lambda t}R_0,\qquad
  R_0:=\norm{e(0)}+\norm{e'(0)}+W(\mu_0,\mu^*)+W(\mu'_0,\mu^*).
\end{equation}
Let $\delta_0:=\norm{x_0-x'_0}+W(\mu_0,\mu'_0)$. Combine this with the coupling estimate of Theorem D(c) on
$[0,T_0]$, with $T_0:=(\Lambda+\lambda)^{-1}\log(CR_0/\delta_0)$.
\begin{itemize}
  \item Both solutions stay in $D$ (proof step~4), so $\norm{x-x_m(t)}<r\le\rho$.
  \item The segment between $x_t$ and $x'_t$ lies in $B(x_m(t),\rho)$, where $f=\fhat^*$, and $\Proj=\mathrm{id}$.
  \item Hence Theorem D(c) holds with $M_x:=k_e+2a_{\max}J_c$, with $L_x$ and $L_\vartheta$ computed for $\bar E=r$
    and an inactive projection, and with $L_2=\lambda_G^{1/2}L_1$. This $\Lambda$ depends neither on $T_0$ nor on $f$
    outside the tube.
\end{itemize}
This gives uniform-in-time Hölder continuity in the basin:
\begin{equation}\label{eq:thmE:holder}
  \sup_{t\ge0}\bigl(\norm{x_t-x'_t}+W(\mu_t,\mu'_t)\bigr)\le(CR_0)^{\Lambda/(\Lambda+\lambda)}\,
  \delta_0^{\lambda/(\Lambda+\lambda)}.
\end{equation}
\end{thm}
\status{Proved, through Theorem E (self-contained since Lemma E.0). Not checked in the repository.}

\begin{proof}
$\norm{x_t-x'_t}=\norm{e-e'}\le\norm{e}+\norm{e'}$, and $W(\mu_t,\mu'_t)\le W(\mu_t,\mu^*)+W(\mu'_t,\mu^*)$. Bound
each bracket by Theorem E. On $[0,T_0]$ use $e^{\Lambda T_0}\delta_0$; on $[T_0,\infty)$ use $CR_0e^{-\lambda T_0}$.
Note that $\delta_0\le R_0\le CR_0$.
\end{proof}

\begin{rmk}{Remark}{E.3}[the collapse point]
Take the $G^{-1}$-closest pair $(j,k)$. Set $e=0$ and $\vartheta_j=\vartheta_k=(\vartheta^*_j+\vartheta^*_k)/2$,
with all other particles ideal.
\begin{itemize}
  \item At this point $\VOT=d^2/8$, and there are two optimal assignments.
  \item By exchange symmetry, $\vartheta_j\equiv\vartheta_k$ for all $t$, so $W(\mu_t,\mu^*)\ge d/2$ forever.
\end{itemize}
The uniqueness argument of step~5 shows $c_0<d^2/8$ automatically, and $d^2/8$ is sharp along this direction.
\end{rmk}
